% concatenated from qaccelerateV238.tex
\documentclass[12pt]{article}

\usepackage{amssymb}%%%%%%%%%%%%%%%
\usepackage{amsmath}%%%%%%%%%%%%%%%
\usepackage{amsthm}%%%%%%%%%%%%%%%

\allowdisplaybreaks%%%%%%%%%%%%%%%

{\theoremstyle{plain}
 \newtheorem{theorem}{Theorem}
}
{\theoremstyle{definition}
\newtheorem{example}{Example}
}

\begin{document}

\begin{center}
{\Large The construction of $q$-analogues via $_3\phi_2$-series and $q$-difference equations}

 \ 

{\textsc{John M. Campbell}}

\end{center}

\begin{abstract}
 We apply the {\tt EKHAD}-normalization method given in our recent work to obtain, via the $q$-version of Zeilberger's algorithm,
 $q$-WZ pairs $(F, G)$ such that $\sum_{k = 0}^{\infty} F(0, k)$ may be expressed as a basic hypergeometric series of the form 
 ${}_{3}\phi_2$ with multiple free parameters, and in such a way so that $\sum_{k=0}^{\infty} F(0, k) = \sum_{n=0}^{\infty} G(n, 0)$. In 
 contrast to how previous applications of {\tt EKHAD}-normalization relied on $q$-analogues for specific WZ pairs introduced by 
 Guillera, our multiparameter approach provides a broad framework in the construction of $q$-analogues for accelerated series for 
 universal constants such as $\pi$. We apply this multiparameter version of {\tt EKHAD}-normalization to obtain and prove new 
 $q$-analogues for accelerated hypergeometric series attributed to many authors, including (alphabetically) Adamchik and Wagon, 
 Ap\'{e}ry, Chu, Chu and Zhang, Fabry, Guillera, Ramanujan, and Zeilberger. 
\end{abstract}

\vspace{0.1in}

\noindent {\footnotesize \emph{MSC:} 33F10, 39A13, 33D15, 33D05}

\vspace{0.1in}

\noindent {\footnotesize \emph{Keywords:} $q$-analogue, hypergeometric series, series acceleration, $q$-difference equation, 
 difference equation, closed form, basic hypergeometric series,  Ramanujan-type series, $q$-shifted factorial}

\section{Introduction}
 The application of $q$-series is of importance within many areas of mathematics, with reference to the text by Gasper and Rahman on 
 basic hypergeometric series \cite{GasperRahman2004} and many subsequent and related research works. As in Gasper and Rahman's 
 monograph \cite[p.\ 36]{GasperRahman2004}, we highlight notable research works in the history of the application of $q$-difference 
 equations, as in the works (chronologically listed) of Jackson \cite{Jackson1910q}, Carmichael \cite{Carmichael1912}, Adams 
 \cite{Adams1931}, Starcher~\cite{Starcher1931}, Trjitzinsky \cite{Trjitzinsky1933}, and Andrews \cite{Andrews1968}. Since
 Wilf--Zeilberger theory marks a paradigm shift \cite{PetkovsekWilfZeilberger1996} given by the use telescoping methods in the field of 
 symbolic computation and beyond \cite{Zeilberger1991}, this motivates new developments and applications of $q$-difference 
 equations via the $q$-Zeilberger algorithm \cite[\S6]{PetkovsekWilfZeilberger1996}. 

 The famous series for $\frac{1}{\pi}$ introduced by Ramanujan in 1914 \cite[pp.\ 352--354]{Berndt1994} \cite{Ramanujan1914} may be 
 seen as prototypical instances of fast converging, hypergeometric series for $\pi$. Given that Ramanujan's series for $\frac{1}{\pi}$ were 
 introduced in 1914, and considering the importance of Ramanujan's series within number theory and other areas, it seems remarkable 
 that the first known $q$-analogues for any of Ramanujan's series were introduced as late as 2018, as in the work of Guo and 
 Liu~\cite{GuoLiu2018}, and as noted by Guo and Zudilin~\cite{GuoZudilin2018}. 
 The methodologies applied in our current paper in the derivation of 
 $q$-analogues related to Ramanujan's summations are inspired by $q$-analogues subsequent to and related to the Guo--Liu 
 $q$-analogues~\cite{ChenChu2021,ChenChu2023,Guillera2018,Guo2020,GuoZudilin2018,HouKrattenthalerSun2019,HouSun2021,
Sun2019,WangZhong2023,Wei2023,Wei2020,WeiRuan2024}. 

 The importance of $q$-analogues for classical, multiparameter, and hypergeometric summation identities is reflected in the above 
 referenced Gasper--Rahman monograph \cite{GasperRahman2004}, with regard to the $q$-Clausen formulas 
 \cite[\S8.8]{GasperRahman2004}, the $q$-Dixon formula \cite[p.\ 44]{GasperRahman2004}, the $q$-Dougall formula~\cite[p.\ 
 44]{GasperRahman2004}, the $q$-Gauss formula \cite[p.\ 14]{GasperRahman2004}, the $q$-Karlsson--Minton formulas 
 \cite[\S1.9]{GasperRahman2004}, the $q$-Kummer formula \cite[p.\ 18]{GasperRahman2004}, the $q$-Saalsch\"{u}tz formula 
 \cite[\S1.7]{GasperRahman2004}, the $q$-Vandermonde formula \cite[p.\ 14]{GasperRahman2004}, the $q$-Watson formula
 \cite[p.\ 61]{GasperRahman2004}, the $q$-Whipple formula~\cite[p.\ 61]{GasperRahman2004}, and many further $q$-analogues. 
 This gives weight to our method for generating $q$-analogues for known hypergeometric evaluations and identities using 
 multiparameter ${}_{3}\phi_{2}$-series. As emphasized via the motivating results in Section \ref{sectionMotivating}, our method can be 
 applied in a very versatile way to obtain new $q$-analogues of previously published results on hypergeometric series. 

 The techniques in this paper build upon and extend our previous use of what we referred to as \emph{{\tt EKHAD}-normalization} 
 \cite{Campbell2026}. Whereas this previous work relied on applying {\tt EKHAD}-normalization using \emph{specific} WZ pairs previously 
 introduced by Guillera \cite{Guillera2006,Guillera2002}, we have formulated a much more general way of applying 
 {\tt EKHAD}-normalization using multiparameter input functions, and this may be seen by analogy with Au's method of WZ seeds 
 \cite{Au2025} and $q$-WZ seeds \cite{Auunpublished} along with our previous work 
 \cite{Campbell2025Hypergeometric,Campbellunpublished,CampbellLevrie2024} related
 to an acceleration method due to Wilf \cite{Wilf1999}. 

\section{Preliminaries and background}
 The $\Gamma$-function may be seen as among the most important special functions, 
 with reference to Rainville's textbook on special functions \cite[\S2]{Rainville1960}, 
 and may be defined 
 via an Euler integral, with 
\begin{equation}\label{displayintegral}
 \Gamma(x) = \int_{0}^{\infty} t^{x-1} e^{-t} \, dt 
\end{equation} 
 for $\Re(x) > 0$. As a notational shorthand, we may write 
\begin{equation}\label{simplifyGamma}
 \Gamma\!\left[ \begin{matrix} \alpha, \beta, \ldots, \gamma \vspace{1mm} \\ 
 A, B, \ldots, C \end{matrix} \right] = \frac{ \Gamma(\alpha) \Gamma(\beta) 
 \cdots \Gamma(\gamma) }{ \Gamma(A) \Gamma(B) \cdots \Gamma(C)}. 
\end{equation}
 The \emph{Pochhammer symbol} may then be defined so that 
\begin{equation}\label{origPoch}
 \left( a \right)_{n} = 
 \Gamma\!\left[ \begin{matrix} a + n \vspace{1mm} \\ 
 a \end{matrix} \right], 
\end{equation}
 and, by analogy with \eqref{simplifyGamma}, we write 
\begin{equation*}
 \left[ \begin{matrix} \alpha, \beta, \ldots, \gamma \vspace{1mm} \\ 
 A, B, \ldots, C \end{matrix} \right]_{k} = \frac{ (\alpha)_{k} (\beta)_{k} 
 \cdots (\gamma)_{k} }{ (A)_{k} (B)_{k} \cdots (C)_{k}}. 
\end{equation*}
 We may define \emph{generalized hypergeometric series} so that 
\begin{equation}\label{pFq}
 {}_{r}F_{s}\!\!\left[ \begin{matrix} 
 a_{1}, a_{2}, \ldots, a_{r} \vspace{1mm}\\ 
 b_{1}, b_{2}, \ldots, b_{s} \end{matrix} \ \Bigg| \ x 
 \right] = \sum_{n=0}^{\infty} 
 \left[ \begin{matrix} a_{1}, a_{2}, \ldots, a_{r} \vspace{1mm} \\ 
 b_{1}, b_{2}, \ldots, b_{s} \end{matrix} \right]_{n} \frac{x^{n}}{n!}, 
\end{equation}
 referring to Bailey's classic text on such series for details and background material \cite{Bailey1964}. We refer to the value $x$ in 
 \eqref{pFq} as the \emph{convergence rate} of either side of \eqref{pFq}. 

 The \emph{$q$-shifted factorial}, which is also referred to as the \emph{$q$-Pochhammer 
 symbol}, is such that 
\begin{equation}\label{qdiscreteshift}
 (a;q)_{n} = \begin{cases} 
 1 & n = 0, \\
 (1-a)(1-aq) \cdots (1-aq^{n-1}) & n \in \mathbb{N}. 
 \end{cases}
\end{equation}
 The definition in \eqref{qdiscreteshift} is also extended so that $$ (a;q)_{\infty} = \prod_{k=0}^{\infty} (1-aq^k) $$ for $q$ such that $|q| < 
 1$, and this allows us to further extend the definition in \eqref{qdiscreteshift} to 
 non-integer values of $n$ according to the relation 
 $$ (a;q)_{n} = \frac{(a;q)_{\infty}}{(aq^n;q)_{\infty}}, $$
 and we make use of the notational shorthand such that 
\begin{equation*}
 \left[ \begin{matrix} \alpha, \beta, \ldots, \gamma \vspace{1mm} \\ 
 A, B, \ldots, C \end{matrix} \, \Bigg| \, q \right]_{n} 
 = \frac{ \left( \alpha; q \right)_{n} 
 \left( \beta; q \right)_{n} \cdots \left( \gamma; q \right)_{n} }{ 
 \left( A; q \right)_{n} 
 \left( B; q \right)_{n} \cdots \left( C; q \right)_{n}}. 
\end{equation*}
 The $q$-Pochhammer symbol gives us a $q$-analogue of the function defined in \eqref{origPoch}
 in the sense that 
 $$ \lim_{q \to 1} \frac{ \left( q^{a};q \right)_{n} }{(1-q)^{n}} = (a)_{n}. $$
 Similarly, the $q$-Gamma function may be defined so that
 $$ \Gamma_{q}(x) = (1-q)^{1-x} \left[ \begin{matrix} q \vspace{1mm} \\ 
 q^{x} \end{matrix} \, \Bigg| \, q \right]_{\infty}, $$
 for $|q| < 1$, giving us a $q$-analogue of \eqref{displayintegral}, with 
 $ \lim_{q \to 1} \Gamma_{q}(x) = \Gamma(x)$, 
 and we adopt a notational shorthand for combinations of $q$-Gamma expressions
 by analogy with \eqref{simplifyGamma}. 
 Also, the \emph{$q$-bracket} 
 is defined so that $[n]_{q} = \frac{1-q^n}{1-q}$, 
 noting that $\lim_{q \to 1} [n]_{q} = n$. 
 By then defining the $q$-factorial
 so that $[n]_{q}! = [1]_{q} [2]_{q} \cdots [n]_{q}$, we may define \emph{Gaussian binomial coefficients}
 so that 
 $$ \left[ \begin{matrix} 
 n \\ 
 k 
 \end{matrix} \right]_{q} = \frac{ [n]_{q}! }{ [k]_{q}! [n-k]_{q}! } $$ for $ k \leq n$. 

 \emph{Unilateral basic hypergeometric series} provide a $q$-analogue of \eqref{pFq} and may be defined so that 
\begin{multline*}
 {}_{j}\phi_{k}\!\!\left[ \begin{matrix} 
 a_{1}, a_{2}, \ldots, a_{j} \vspace{1mm}\\ 
 b_{1}, b_{2}, \ldots, b_{k} \end{matrix} \ \Bigg| \ q; z 
 \right] \\ 
 = \sum_{n=0}^{\infty} \frac{z^n}{ (q;q)_{n} } 
 \left[ \begin{matrix} 
 a_{1}, a_{2}, \ldots, a_{j} \vspace{1mm} \\ 
 b_{1}, b_{2}, \ldots, b_{k} \end{matrix} \, \Bigg| \, q \right]_{n} 
 \left( (-1)^{n} q^{\binom{n}{2}} \right)^{1 + k - j}. 
\end{multline*}
 Basic hypergeometric series of the form ${}_{3}\phi_{2}$ are to play a key role in our work. 
 Examples of especially notable ${}_{3}\phi_{2}$-identities include
 the relation 
$$ {}_{3}\phi_{2}\!\!\left[ \begin{matrix} 
 q^{-n}, a, b \vspace{1mm}\\ 
 c, 	 \frac{a b q^{1-n}}{c} \end{matrix} \ \Bigg| \ q; 1 
 \right] 
 = \left[ \begin{matrix} 
 \frac{c}{a}, \frac{c}{b} \vspace{1mm} \\ 
 c, \frac{c}{ab} \end{matrix} \, \Bigg| \, q \right]_{n} 
 $$ for integers $n \geq 0$ attributed to Jackson \cite{Jackson1910Transformations}. 
 Notable instances of ${}_{3}\phi_{2}$-transforms include 
 the $q$-analogue
\begin{multline*} 
 {}_{3}\phi_{2}\!\!\left[ \begin{matrix} 
 a, b, c \vspace{1mm}\\ 
 d, e \end{matrix} \ \Bigg| \ q; q^{d+e-a-b-c} 
 \right] 
 = \\ \Gamma_{q}\!\left[ \begin{matrix} e, d+e-a-b-c \vspace{1mm} \\ 
 e-a, d + e - b - c \end{matrix} \right] 
 {}_{3}\phi_{2}\!\!\left[ \begin{matrix} 
 a, d-b, d-c \vspace{1mm}\\ 
 d, d+e-b-c \end{matrix} \ \Bigg| \ q; q^{e-a} 
 \right] 
\end{multline*} 
 of Thomae's transformation \cite[p. 200]{GasperRahman2004}, referring to Gasper and Rahman's text 
 for background and for many related transforms. 

\subsection{{\tt EKHAD}-normalization}\label{sectionEKHAD}
 A bivariate function $F(n, k)$ is \emph{hypergeometric} if $\frac{F(n+1,k)}{F(n,k)}$ and $\frac{F(n,k+1)}{F(n,k)}$ are both rational functions, 
 and the coefficients are taken to be in $\mathbb{Q}$ in the context of WZ theory. For a rational function $R(n, k)$, write $G(n, k) = R(n, 
 k) F(n, k)$, and suppose that $F$ and $G$ satisfy a difference equation of the form 
\begin{equation}\label{fundamentalWZ}
 F(n+1, k) - F(n, k) = G(n, k+1) - G(n, k). 
\end{equation}
 This may be seen as the fundamental difference equation 
 in WZ theory, with reference to the classic text by 
 Petkov\v sek et al.\ on WZ theory \cite{PetkovsekWilfZeilberger1996}. 
 While a given function $F(n, k)$ may not be such that there exists a companion $G(n, k)$
 of the desired form satisfying \eqref{fundamentalWZ}, 
 \emph{Zeilberger's algorithm} \cite[\S6]{PetkovsekWilfZeilberger1996} 
 may be thought of as extending the WZ difference equation in \eqref{fundamentalWZ}
 in such a way so that this algorithm
 is applicable to \emph{any} hypergeometric input function $F(n, k)$, 
 referring to the same text by Petkov\v sek et al.\ for details. 
 Our techniques rely on the $q$-version of Zeilberger's algorithm, 
 again referring to the appropriate text by Petkov\v sek et al.\ for details. 

 A \emph{first-order} difference equation
 produced by Zeilberger's algorithm (resp.\ the $q$-Zeilberger algorithm)
 with $F(n, k)$ as a given input function, assuming that $F(n, k)$
 is such that a difference equation of the following form exists, 
 is such that 
\begin{equation}
 p_{1}(n) F(n+1, k) - p_{2}(n) F(n, k) = G(n, k+1) - G(n, k) 
\end{equation}
 for (single-variable) polynomials 
 $p_{1}(n)$ and $p_{2}(n)$ (resp.\ $q$-analogues). 
 Following our past work \cite{Campbell2026}, by ``normalizing'' 
 the input function $F(n, k)$ by setting 
\begin{equation}\label{displaynormal}
 \overline{F}(n, k) = (-1)^{n} \left( \prod_{i=j}^{n-1} \frac{p_{1}(i)}{p_{2}(i)} \right) F(n, k) 
\end{equation}
 for minimal $j$ such that $p_2$ is nonvanishing, 
 typically with $j \in \{ 0, 1 \}$, 
 and by applying Zeilberger's algorithm to the normalized version of $F$ in \eqref{displaynormal}, 
 we obtain a companion function
 $\overline{G}(n, k)$ such that $\overline{F}$ and $\overline{G}$ satisfy the desired WZ difference equation, 
 and this may be verified inductively. 
 As in our past work \cite{Campbell2026}, we refer to the normalization process indicated above
 as \emph{{\tt EKHAD}-normalization}. 


 Our past work \cite{CampbellLevrie2024} concerned multiparameter versions of an acceleration method due to Wilf 
 \cite{Wilf1999}.     
 In   this     past work \cite{CampbellLevrie2024}, by applying Zeilberger's algorithm
 to certain input functions such as 
 \begin{equation}\label{freelowerquarter}
 F(n, k) := \left[ \begin{matrix} 
 a, b \vspace{1mm} \\ 
 n, n
 \end{matrix} \right]_{k}, 
\end{equation}
 this produced inhomogeneous, first-order difference equations that were \emph{not}
 normalized, in accordance with Wilf's method \cite{Wilf1999}. 
 Due to certain conver-gence-related issues, it appears that this approach 
 would not work, in general, be applicable using $q$-analogues of input function as in 
 \eqref{freelowerquarter}, 
 and this leads us to apply {\tt EKHAD}-normalization to $q$-analogues of 
 hypergeometric expressions as in \eqref{freelowerquarter}, 
 and, instead of using Wilf's acceleration method (which relies on inhomogeneous difference equations, 
 in contrast to \eqref{fundamentalWZ}), we instead consider a telescoping argument 
 based on WZ difference equations. 

 For a $q$-analogue $\mathcal{F}(n, k; q)$ of a hypergeometric input function as in 
 \eqref{freelowerquarter}, 
 suppose that $\mathcal{F}(n, k; q)$ satisfies a first-order difference equation according
 to the $q$-Zeilberger algorithm.
 By then applying {\tt EKHAD}-normalization to $\mathcal{F}(n, k; q)$, 
 yielding a $q$-WZ pair $(\overline{\mathcal{F}}, \overline{\mathcal{G}})$, a telescoping phenomenon gives us that 
\begin{equation*}
 \overline{\mathcal{F}}(m+1, k; q) - \overline{\mathcal{F}}(0, k; q) = \sum_{n=0}^{m} \overline{\mathcal{G}}(n, k+1; q) - \sum_{n=0}^{m} \overline{\mathcal{G}}(n, k; q). 
\end{equation*}
 Depending on the parameters, 
 suppose that $\lim_{m \to \infty} \overline{\mathcal{F}}(n+1,k;q)$ vanishes, with 
\begin{equation}\label{meta1}
 - \overline{\mathcal{F}}(0, k; q) = \sum_{n=0}^{\infty} \overline{\mathcal{G}}(n, k+1; q) - \sum_{n=0}^{\infty} \overline{\mathcal{G}}(n, k; q) 
\end{equation}
 and 
\begin{equation}\label{meta2}
 - \overline{\mathcal{F}}(0, k+1; q) = \sum_{n=0}^{\infty} \overline{\mathcal{G}}(n, k+2; q) - \sum_{n=0}^{\infty} \overline{\mathcal{G}}(n, k+1; q) 
\end{equation}
 and 
\begin{equation}\label{meta3}
 - \overline{\mathcal{F}}(0, k+2; q) = 
 \sum_{n=0}^{\infty} \overline{\mathcal{G}}(n, k+3; q) - \sum_{n=0}^{\infty} \overline{\mathcal{G}}(n, k+2; q), 
\end{equation}
 and so forth. Ideally, and again 
 depending on the input parameters, 
 we would apply a telescoping argument 
 suggested in 
 \eqref{meta1}--\eqref{meta3} together with 
 an analytic argument to prove the permissibility of an interchange of limiting operations of the form 
 $$\lim_{m \to \infty} \sum_{n=0}^{\infty} \overline{\mathcal{G}}(n,k+m;q) = 
 \sum_{n=0}^{\infty} \lim_{m \to \infty} \overline{\mathcal{G}}(n,k+m;q),$$ 
 to prove that 
\begin{equation}\label{FkGn}
 \sum_{k=0}^{\infty} \overline{\mathcal{F}}(0, k;q) = \sum_{n=0}^{\infty} \overline{\mathcal{G}}(n,0;q). 
\end{equation}
 We apply this approach 
 to produce the motivating results highlighted in Section \ref{sectionMotivating}
 along with many further results. 
 Through extensive computer searches for combinations of input parameters such that 
 the ``non-$q$'' version of \eqref{FkGn}
 admits a closed form, 
 this provides a versatile way of
 generating new $q$-analogues, as demonstrated in Section 
 \ref{sectionMotivating} below. 

\subsection{Ramanujan-type formulas and $\Gamma$-values} 
 Ramanujan's $17$ formulas for $\frac{1}{\pi}$ \cite[pp.\ 352--354]{Berndt1994} \cite{Ramanujan1914} include 
\begin{equation}\label{Ramanujanquarter}
 \frac{4}{\pi} = 
 \sum_{n = 0}^{\infty} 
 \left( \frac{1}{4} \right)^{n} \left[ \begin{matrix} 
 \frac{1}{2}, \frac{1}{2}, \frac{1}{2} 
 \vspace{1mm} \\ 
 1, 1, 1
 \end{matrix} \right]_{n} (6 n + 1), 
 \end{equation}
 and one of the main results in our paper is given by a new 
 $q$-analogue of \eqref{Ramanujanquarter} 
 that we highlight in Section \ref{subRamanujan}
 and that we prove as a consequence of Theorem \ref{theoremquarter} below. 
 Further formulas, out of Ramanujan's $17$ formulas for $\frac{1}{\pi}$, include
 $$ \frac{27}{4\pi} = \sum_{n = 0}^{\infty} 
 \left( \frac{2}{27} \right)^{n} \left[ \begin{matrix} 
 \frac{1}{3}, \frac{1}{2}, \frac{2}{3} 
 \vspace{1mm} \\ 
 1, 1, 1
 \end{matrix} \right]_{n} (15 n + 2) 
 $$ and 
 $$ \frac{72}{\pi} = \sum_{n = 0}^{\infty} 
 \left( -\frac{1}{324} \right)^{n} \left[ \begin{matrix} 
 \frac{1}{4}, \frac{1}{2}, \frac{3}{4} 
 \vspace{1mm} \\ 
 1, 1, 1
 \end{matrix} \right]_{n} (260 n + 23). $$

 The importance of Ramanujan's $\pi$-formulas in the history of mathematics, especially in regard to subjects such as computational 
 number theory, is made clear through the expository work 
 due to Baruah et al.\ \cite{BaruahBerndtChan2009}
 along with the classic monograph on $\pi$ and the AGM \cite{BorweinBorwein1987}. 
 This motivates the concept of a 
 \emph{Ramanujan-type formula}, which, for our purposes, refers to 
 an evaluation, in closed form or in terms of constants such as $\Gamma$-values with rational arguments, 
 for a series of the form 
\begin{equation}\label{20260213945PM12AS}
 \sum_{n=0}^{\infty} z^n \frac{ (a_{1})_{n} (a_{2})_{n} (a_{3})_{n} }{ 
 (b_{1})_{n} (b_{2})_{n} (b_{3})_{n} } (c_{1} n + c_2), 
\end{equation}
 with $a_i$ and $b_i$ and $c_i$ being rational, 
 and with the condition that 
 the difference between the arguments of any lower 
 and upper Pochhammer symbols within the summand in \eqref{20260213945PM12AS} is not an integer. 
 Many of our $q$-analogues introduced in this paper
 are $q$-analogues of known Ramanujan-type formulas. 
 Moreover, our techniques have the added benefit of producing
 many further Ramanujan-type formulas that have not previously been known. 

 While a number of our past works have concerned closed-form evaluations for series satisfying
 the given conditions for \eqref{20260213945PM12AS}, 
 there is relatively little known when it comes to the expression of 
 series satisfying the conditions for 
 \eqref{20260213945PM12AS} with $\Gamma$-values with rational arguments. 
 We apply our techniques given in this paper 
 to produce many new evaluations of this form, 
 and this is motivated by past and remarkable evaluations of this form 
 due to Chen and Chu \cite{ChenChu2024}, including 
\begin{align*}
 \frac{\pi \Gamma^{2}\left( \frac{1}{3} \right)}{ 6 \Gamma^{2}\left( \frac{5}{6} 
 \right) } & = \sum_{n = 0}^{\infty} 
 \left( -\frac{1}{27} \right)^{n} \left[ \begin{matrix} 
 \frac{1}{2}, \frac{2}{3}, 1 
 \vspace{1mm} \\ 
 \frac{4}{3}, \frac{4}{3}, \frac{4}{3}
 \end{matrix} \right]_{n} (7n+3), \\ 
 \frac{2\pi^2}{\Gamma^{3}\left( \frac{2}{3} \right)} & = \sum_{n = 0}^{\infty} 
 \left( -\frac{1}{27} \right)^{n} \left[ \begin{matrix} 
 \frac{1}{6}, \frac{2}{3}, \frac{2}{3} 
 \vspace{1mm} \\ 
 1, \frac{7}{6}, \frac{4}{3}
 \end{matrix} \right]_{n} (21n+8), \ \text{and} \\ 
 \frac{2\pi^2}{\Gamma^{3}\left( \frac{1}{3} \right)} & = \sum_{n = 0}^{\infty} 
 \left( -\frac{1}{27} \right)^{n} \left[ \begin{matrix} 
 -\frac{1}{6}, \frac{1}{3}, \frac{1}{3} 
 \vspace{1mm} \\ 
 \frac{2}{3}, \frac{5}{6}, 1
 \end{matrix} \right]_{n} (21n+1). 
\end{align*}
 Our $q$-analogues of Ramanujan-type formulas for $\Gamma$-values with rational 
 arguments are also inspired by $q$-analogues of this form due to 
 Chen and Chu~\cite{ChenChu2023}, 
 with reference to their applications of 
 Carlitz inversions of
 the $q$-Pfaff-Saalsch\"{u}tz theorem. 
 This approach led Chen and Chu to introduce evaluations such as 
 $$ \frac{\Gamma^{2}\left( \frac{1}{4} \right)}{8 \Gamma^{2}\left( \frac{3}{4} \right) } 
 = \sum_{n = 0}^{\infty} 
 \left( \frac{1}{4} \right)^{n} \left[ \begin{matrix} 
 \frac{1}{2}, \frac{1}{2}, \frac{1}{2} 
 \vspace{1mm} \\ 
 1, \frac{5}{4}, \frac{5}{4}
 \end{matrix} \right]_{n} (3n+1). $$

\section{Motivating results}\label{sectionMotivating}
 We highlight, in the current section, new results obtained from our multiparameter {\tt EKHAD}-normalization method that provide 
 $q$-analogues of notable summation identities from a number of different authors. 

 Ramanujan's formulas for $\frac{1}{\pi}$ are groundbreaking in terms of the history associated with the computation of $\pi$, with 
 reference to relevant texts related to this history 
 \cite{ArndtHaenel2001,BerggrenBorweinBorwein2004,BorweinBorwein1987}. 
 Since our work concerns $q$-analogues of fast converging series for universal constants such as $\pi$, 
 the groundbreaking nature of the 
 Guo--Liu $q$-analogues introduced in 2018 
 leads us toward our new $q$-analogue of the simplest or slowest converging
 out of Ramanujan's 17 series for $\frac{1}{\pi}$. 
 
\subsection{Ramanujan's series of convergence rate $\frac{1}{4}$}\label{subRamanujan}
 In Section \ref{subquarter}, we prove the relation 
\begin{multline*}
 q - 1 + \frac{q-1}{(q+1)^2} \sum_{n=0}^{\infty} 
 q^{2n} \left[ \begin{matrix} 
 q, q \vspace{1mm} \\ 
 q^4, q^4 \end{matrix} \, \Bigg| \, q^{2} \right]_{n} = \\ \sum_{n=0}^{\infty} q^{2n} 
 \left[ \begin{matrix} 
 q, q, q, q^3 \vspace{1mm} \\ 
 q^2, q^2 \end{matrix} \, \Bigg| \, q^{2} \right]_{n} 
 \frac{q^{4n+1} + q^{2n} - 2}{ (q^4;q^4)_{n} (q^6;q^4)_{n}} 
\end{multline*}
 for $|q| > 1$. This provides a new $q$-analogue of 
 the Ramanujan relation in \eqref{Ramanujanquarter}, with regard to the survey below. 

 The above $q$-analogue of \eqref{Ramanujanquarter} is not equivalent to 
 the previously known $q$-analogues of \eqref{Ramanujanquarter}, 
 namely, the Guo--Liu $q$-analogue of \eqref{Ramanujanquarter} \cite{GuoLiu2018} such that 
$$ \sum_{n=0}^{\infty} 
 q^{n^2} [6n+1] \left[ \begin{matrix} 
 q^2 \vspace{1mm} \\ 
 q^4, q^4, q^4 \end{matrix} \, \Bigg| \, q^{4} \right]_{n} (q;q^2)_{n}^{2} 
 = (1+q) 
 \left[ \begin{matrix} 
 q^{2}, q^{6} \vspace{1mm} \\ 
 q^4, q^4 \end{matrix} \, \Bigg| \, q^{4} \right]_{\infty}, $$
 the Guo $q$-analogue of \eqref{Ramanujanquarter} \cite{Guo2020} such that 
\begin{multline*}
 \sum_{n=0}^{\infty} 
 q^{2n^2} [6n+1] \left[ \begin{matrix} 
 q, q, q, q \vspace{1mm} \\ 
 q^2, q^2 \end{matrix} \, \Bigg| \, q^{2} \right]_{n} \frac{1}{(q^2;q^2)_{2n}} \times \\ 
 \left( [4n+1] - \frac{q^{4n+1}[2n+1]^2}{[4n+2]} \right) = 
 \left[ \begin{matrix} 
 q, q^{3} \vspace{1mm} \\ 
 q^2, q^2 \end{matrix} \, \Bigg| \, q^{2} \right]_{\infty}, 
\end{multline*}
 our previous $q$-analogue of \eqref{Ramanujanquarter} \cite{Campbell2026} such that 
\begin{multline*}
 \sum_{n=0}^{\infty} (-1)^n 
 q^{n^2} \frac{ (1-q^{ -1 + 2n}) (2 q^{1+4n} - q^{1+2n} - 1 ) }{1 + q^{1+2n}} \\ 
 \times \left[ \begin{matrix} \frac{1}{q}, q, q, q, q \vspace{1mm} \\ 
 q^2, q^2, q^2, q^{1-2n}, q^{2+2n} \end{matrix} \, \Bigg| \, q^2 \right]_{n}
 = 
 \frac{(1-q)^2 (1+q)}{q} \left[ \begin{matrix} q, q \vspace{1mm} \\ 
 q^2, q^3 \end{matrix} \, \Bigg| \, q^2 \right]_{\frac{1}{2}}, 
\end{multline*}
 our previous $q$-analogue of \eqref{Ramanujanquarter} 
 \cite{Campbell2026} such that
\begin{multline*}
 \sum_{n=0}^{\infty} 
 q^{2 n^2} \frac{ (q - q^{2n}) (2q^{1+4n} - q^{1+2n} - 1 ) }{ (1+q^{2n}) (1 + q^{1+2n}) } \\
 \times \left[ \begin{matrix} \frac{1}{q}, q, q \vspace{1mm} \\ 
 -1, -q, q^2, q^2, q^2 \end{matrix} \, \Bigg| \, q^2 \right]_{n} 
 = 
 \frac{(1-q)^2 (1+q) }{2} \left[ \begin{matrix} q, q \vspace{1mm} \\ 
 q^2, q^3 \end{matrix} \, \Bigg| \, q^2 \right]_{\frac{1}{2}}, 
\end{multline*}
 Au's $q$-analogue of \eqref{Ramanujanquarter}
 such that 
$$ \sum_{n=0}^{\infty} 
 q^{2n^2} \frac{ \left( q;q^2 \right)_{n}^{4} }{ \left( q^2;q^2 \right)_{2n} 
 \left( q^2;q^2 \right)_{n}^{2} } \frac{1+q^{2n+1}-2q^{4n+1}}{1+q^{2n+1}} 
 = \left[ \begin{matrix} q, q \vspace{1mm} \\ 
 q^2, q^2 \end{matrix} \, \Bigg| \, q^2 \right]_{\infty} $$
 Au's $q$-analogue of \eqref{Ramanujanquarter}
 such that 
$$ \sum_{n=0}^{\infty} 
 q^{ n^2} \frac{ \left( q;q^2 \right)_{n}^{2} \left( q^2;q^4 \right)_{n} }{ 
 \left( q^4;q^4 \right)_{n}^{3} } (1-q^{6n+1}) 
 = \left[ \begin{matrix} q^2, q^2 \vspace{1mm} \\ 
 q^4, q^4 \end{matrix} \, \Bigg| \, q^4 \right]_{\infty}, $$
 and Au's $q$-analogue of \eqref{Ramanujanquarter}
 such that 
$$ \sum_{n=0}^{\infty} 
 q^{ 2 n} \frac{ \left( q;q^2 \right)_{n}^{4} }{ 
 \left( q^2;q^2 \right)_{2n} \left( q^2;q^2 \right)_{n}^{2} } \frac{2-q^{2n} - q^{4n+1}}{1+q^{2n+1}} 
 = \left[ \begin{matrix} q, q, q, q \vspace{1mm} \\ 
 q^2, q^2, q^2 \end{matrix} \, \Bigg| \, q^2 \right]_{\infty} \sum_{n=1}^{\infty} \frac{q^{n-1}}{ \left( q;q^2 \right)_{n} }. $$
 The above Guo $q$-analogue of Ramanujan's formula was proved in an alternative way by Chen and Chu 
 through the use of Carlitz inversions of
 the $q$-Pfaff-Saalsch\"{u}tz theorem \cite{ChenChu2023}. 

\subsection{The Fabry--Guillera formula}\label{subFabryGuillera}
 In Section \ref{subquarter}, we prove the relation 
\begin{multline*}
 \frac{(1-q)^3 (1+q) \left(1 + q^2\right)}{q}
 \sum_{k = 0}^{\infty} 
 \frac{q^{2k}}{(1-q^{2k+1})^{2} } = \\ \sum_{n=0}^{\infty} q^{2n} 
 \left[ \begin{matrix} 
 q^{2}, q^{2}, q^{2}, q^{4} \vspace{1mm} \\ 
 q^{3}, q^{3} \end{matrix} \, \Bigg| \, q^{2} \right]_{n} 
 \frac{2-q^{2n+1} - q^{4n+3}}{(q^6;q^4)_{n} (q^8;q^4)_{n}} 
\end{multline*}
 for $|q| > 1$. This provides a $q$-analogue of the Ramanujan-type formula 
\begin{equation}\label{FabryGuillera}
 \frac{\pi^2}{4} = 
 \sum_{n = 0}^{\infty} 
 \left( \frac{1}{4} \right)^{n} \left[ \begin{matrix} 
 1, 1, 1 
 \vspace{1mm} \\ 
 \frac{3}{2}, \frac{3}{2}, \frac{3}{2}
 \end{matrix} \right]_{n} (3 n + 2), 
 \end{equation}
 with regard to the survey below, and 
 with \eqref{FabryGuillera} having originally been due (in an equivalent form) to Fabry in 1911
 \cite[p.\ 129]{Fabry1911} 
 and having subsequently been attributed by many authors Guillera, 
 who independently rediscovered \eqref{FabryGuillera} via the WZ method in 2008 \cite{Guillera2008}. 
 We also introduce a further $q$-analogue of \eqref{FabryGuillera} whereby 
\begin{multline*}
 (1-q) (1-q^2) 
 \sum_{k = 0}^{\infty} 
 \left[ \begin{matrix} 
 q^2 \vspace{1mm} \\ 
 q^{3} \end{matrix} \, \Bigg| \, q^{2} \right]_{k} \frac{q^{2k}}{1-q^{2k+2}} = \\ 
 \sum_{n=0}^{\infty} q^{2n} 
 \frac{ \left( q^2;q^2 \right)_{n}^{4} }{ \left( q^{3};q^{2} \right)_{n}^{2} 
 \left( q^{4};q^{4} \right)_{n} \left( q^{6};q^{4} \right)_{n} }
 \frac{2 - q^{2n+1} - q^{4n+3}}{1 + q^{2n+2}} 
\end{multline*}
 for $|q| > 1$. 

 There have been a number of past research contributions involving a variety of different methods in the derivation of $q$-analogues 
 of equivalent formulations of the Fabry--Guillera formula in \eqref{FabryGuillera}, 
 which is representative of the research interest in the new $q$-analogue shown above
 of the Fabry--Guillera formula. 
 Notably, Hou et al.\ \cite{HouKrattenthalerSun2019} introduced and proved 
 the $q$-analogue
\begin{multline*}
 \frac{1}{2} \sum_{n=0}^{\infty} \frac{q^{2n}}{(1-q^{2n+1})^{2}} = \\ 
 \sum_{n=0}^{\infty} q^{2n(n+1)} 
 \frac{ \left( q^{2};q^2 \right)_{n}^{3} }{ \left( q;q^2 \right)_{n+1}^{3} 
 \left( -1; q \right)_{2n+3} } (1+q^{2n+2} - 2 q^{4n+3})
\end{multline*}
 of the Fabry--Guillera formula, along with the $q$-analogue 
\begin{equation}\label{secondHKS}
 \sum_{n=0}^{\infty} q^{\frac{n(n+1)}{2}} 
 \frac{1-q^{3n+2}}{1-q} \frac{ \left( q;q \right)_{n}^{3} \left( -q;q \right)_{n} }{ 
 \left( q^{3};q^{2} \right)_{n}^{3} } 
 = (1-q)^2 \left[ \begin{matrix} 
 q^{2}, q^{2}, q^{2}, q^{2} \vspace{1mm} \\ 
 q, q, q, q \end{matrix} \, \Bigg| \, q^{2} \right]_{\infty}
\end{equation}
 of the Fabry--Guillera formula. 

 Chen and Chu \cite{ChenChu2021} obtained an equivalent version of \eqref{secondHKS} using the relation 
\begin{multline*}
 \sum_{k=0}^{\infty} 
 \frac{1-q^{3k}a}{1-a} 
 \left[ \begin{matrix} 
 b, d, \frac{qa}{bc} \vspace{1mm} \\ 
 q \end{matrix} \, \Bigg| \, q \right]_{k} 
 \left[ \begin{matrix} 
 a \vspace{1mm} \\ 
 \frac{q^2 a}{b}, \frac{q^2a}{d}, q b d \end{matrix} \, \Bigg| \, q^{2} \right]_{k} q^{\binom{k+1}{2}} 
 = \\ 
 \left[ \begin{matrix} 
 q^2 a, qb, qd, \frac{q^2a}{bd} \vspace{1mm} \\ 
 q, \frac{q^2a}{b}, \frac{q^2 a}{d}, q b d \end{matrix} \, \Bigg| \, q^{2} \right]_{\infty} 
\end{multline*}
 due to Rahman \cite{Rahman1993}. 
 Chu \cite{Chu2018} also obtained an equivalent version of 
 \eqref{secondHKS} 
 using recursions for $q$-analogues of the $\Omega$-sums 
 involved in previous references from Chu et al., 
 as in the work of Chu and Zhang on the acceleration of Dougall's ${}_{5}F_{4}$-summation
 \cite{ChuZhang2014}. 
 Wei \cite{Wei2020} also obtained an equivalent formulation of 
 \eqref{secondHKS} 
 using a terminating $q$-hypergeometric relation due to Chu~\cite{Chu1995} 
 via an inversion technique applied using Jackson's $q$-analogue 
 of the Dougall--Dixon theorem. 
 The variety of published proofs of the $q$-analogue 
 in \eqref{secondHKS} 
 of the Fabry--Guillera formula
 motivate the $q$-analogue of this formula we introduce, 
 via {\tt EKHAD}-normalization. 
 
 Wang and Zhong \cite{WangZhong2023} introduced the $q$-analogue 
\begin{multline*}
 \sum_{n=0}^{\infty} q^{\frac{n^2 + 5n+2}{2}} \frac{ [3n+4]_{q} }{ [2n+3]_{q}^{2} } 
 \frac{ \left( q;q \right)_{n}^{3} \left( -q;q \right)_{n} }{ \left( q^{3};q^{2} \right)_{n}^{3} } 
 = \\ (1-q)^2 \left[ \begin{matrix} 
 q^2, q^2, q^2, q^2 \vspace{1mm} \\ 
 q, q, q, q \end{matrix} \, \Bigg| \, q^{2} \right]_{\infty} - (1+q) 
\end{multline*}
 of an equivalent version of the Fabry--Guillera formula, using a $q$-rational reduction technique. 
 In our previous work \cite{Campbell2025further}, 
 we introduced 
 the $q$-series identity such that 
\begin{multline*} 
 \frac{q^{(k-1)^2}}{(q;q)_{k-1}^2} \sum_{n=0}^{\infty} 
 \frac{q^{n} \left(2 q^k-q^{n+1}-q^{2 n+2}\right)}{1+q^{n+1}} 
 \frac{(q;q)_n^4}{ (q;q)_{2 n+1} (q;q)_{n-k+1}^2} = \\ 
 \frac{1}{(q;q)_{\infty }^4} \sum_{n=0}^{\infty} q^{(k+n)^2} \left((q;q)_{\infty }^3-\left(q^{1-n - 
 k};q\right)_{\infty }^2\right) \left(q^{1+n+k};q\right)_{\infty }^2, 
\end{multline*}
 with the $k = \frac{1}{2}$ providing a $q$-analogue of the Fabry--Guillera formula. 

\subsection{Ap\'{e}ry's series for $\zeta(2)$}\label{subApery}
 We introduce the $q$-analogue $$ (1-q)^2 (1+q) \sum_{n=0}^{\infty} 
 \frac{q^n}{(1-q^{n+1})^2} = \sum_{n=0}^{\infty} q^{n} 
 \frac{ (q;q)_{n}^{2} }{ (q^3;q^2)_{n} (q^4;q^2)_{n} } 
 (2+q^{n+1}) $$
 for $|q| > 1$ 
 of the formula 
\begin{equation}\label{originalApery}
 \frac{\pi^2}{9} = 
 \sum_{k = 0}^{\infty} 
 \left( \frac{1}{4} \right)^{n} \left[ \begin{matrix} 
 1 
 \vspace{1mm} \\ 
 \frac{3}{2} 
 \end{matrix} \right]_{k} \frac{1}{k+1} 
 \end{equation}
 attributed to Ap\'{e}ry, as in the work of Chu and Zhang \cite[Example 74]{ChuZhang2014}. An equivalent version (up to a reindexing) of 
 \eqref{originalApery} played a key step in Ap\'{e}ry's renowed proof \cite{Apery1979,vanderPoorten197879} of the irrationality of 
 $\zeta(2)$ for the Riemann zeta function $\zeta(x) = \sum_{n=1}^{\infty} \frac{1}{n^x}$. This motivates the new $q$-analogue of
 Ap\'{e}ry's formula given in this paper,  along with number-theoretic applications of this $q$-analogue.  Our $q$-analogue of Ap\'{e}ry's  
 series for $\zeta(2)$ is also motivated by Hessami Pilehrood and Hessami Pilehrood's   $q$-analogue  
 \cite{HessamiPilehroodHessamiPilehrood2011} of the Bailey--Borwein--Bradley   identity, which can be thought of as an extension of  
 Ap\'{e}ry's series for $\zeta(2)$ and $\zeta(3)$. 

\subsection{Zeilberger's series of convergence rate $\frac{1}{64}$}\label{sectionZeilberger}
 A remarkable result due to Zeilberger is such that 
\begin{equation}\label{displayZeilberger}
 \frac{4\pi^2}{3} = 
 \sum_{k = 0}^{\infty} 
 \left( \frac{1}{64} \right)^{n} \left[ \begin{matrix} 
 1, 1, 1 
 \vspace{1mm} \\ 
 \frac{3}{2}, \frac{3}{2}, \frac{3}{2} 
 \end{matrix} \right]_{k} (21k + 13) 
 \end{equation}
 and was introduced in 1993 \cite{Zeilberger1993} via the WZ method and was subsequently obtained by Guillera \cite{Guillera2008}, 
 again by the WZ method. We introduce the $q$-analogue 
\begin{multline*}
 -(1-q)^3 \sum_{k=0}^{\infty} \frac{q^{k}}{(1-q^{k+1})^2} = \sum_{n=0}^{\infty} q^{n} 
 \frac{ \left( q;q \right)_{n}^{6} }{ \left( q^{2};q^{2} \right)_{n}^{3} 
 \left( q^{3};q^{2} \right)_{n}^{3} } \times \\ 
 \frac{-3 q^{n+1}+q^{2 n+1}-2 q^{2 n+2}+3 q^{3 n+2}+3 q^{4 n+3}+q^{5 n+4}-3}{\left(q^{n+1}+1\right)^3} 
\end{multline*}
 for $|q| > 1$ of the Zeilberger formula in \eqref{displayZeilberger}. This may be considered in relation to the previously
 known $q$-analogue 
\begin{multline*}
 (1-q)^{2} \sum_{n=1}^{\infty} \frac{ n q^{n} }{1-q^{n}} = \sum_{n=1}^{\infty} \frac{q^{3n^2-3n+1}}{ \left[ \begin{smallmatrix} 
 2 n \\ 
 n 
 \end{smallmatrix} \right]_{q}^{3} [n]_{q}^{3} } 
 \big( (3q^{4n-1} + 6 q^{3n-1} - \\ q^{2n} + 8 q^{2n-1} - 4 q^{n} + 8 q^{n-1} + 1) [n]_{q} - 8 q^{n} \big) 
\end{multline*}
 of \eqref{displayZeilberger} introduced and proved by 
 Hessami Pilehrood and Hessami Pilehrood in 2011 \cite{HessamiPilehroodHessamiPilehrood2011}. 

 \subsection{A $q$-analogue of a BBP-type formula}
 The new $q$-identity 
\begin{multline*}
 q(1 - q^2) 
 \sum_{n=0}^{\infty} q^{2n}
 \left[ \begin{matrix} 
 q \vspace{1mm} \\ 
 q^2 \end{matrix} \, \Bigg| \, q^2 \right]_{n} 
 \frac{1}{1-q^{2n+1}} = \sum_{n=0}^{\infty} (-1)^{n} 
 q^{-n^2-n} \times \\ 
 \frac{ (q;q^2)_{n} (q^2;q^2)_{n} }{(q^4;q^4)_{n} (q^6;q^4)_{n} } 
 \frac{1 - q^{6 n+4}-q^{8 n+3}-q^{8 n+4}+q^{10 n+5}+q^{12 n+6}}{\left(1-q^{4 n+1}\right) \left(1-q^{4 n+3}\right)}. 
\end{multline*}
 for $|q| > 1$ is a $q$-analogue of the BBP-type formula $$ \pi = \sum_{n=0}^{\infty} \left( -\frac{1}{4} \right)^{n} \left( \frac{2}{4n+1} + 
 \frac{2}{4n+2} + \frac{1}{4n+3} \right) $$ attributed to Adamchik and Wagon \cite{AdamchikWagon1996}. This provides a 
 companion to a $q$-analogue due to Chu \cite{Chu2018} of the Adamchik--Wagon BBP-type formula $$ \frac{3\pi}{2} = \sum_{n = 
 0}^{\infty} \left( -\frac{1}{4} \right)^{n} \left( \frac{1}{4n+1} + \frac{2}{4n+2} + \frac{1}{4n+3} \right). $$ Apart from the work of Chu 
 \cite{Chu2018}, there does not seem to be much known about $q$-analogues of BBP-type formulas. 

\section{Multiparameter {\tt EKHAD}-normalization}
 Our $q$-analogues for accelerated hypergeometric series are organized as special cases of the Theorems presented below. 

\subsection{$q$-analogues for the convergence rate $\frac{1}{4}$}\label{subquarter}
 By analogy with how Wilf's acceleration method \cite{Wilf1999} was applied in the first author's previous work by setting $F(n, k)$ as in 
 for free parameters $a$ and $b$ and by applying Zeilberger's algorithm to \eqref{freelowerquarter} to produce a first-order difference, 
 inhomogeneous difference equation, we require a homogeneous and $q$-WZ difference equation to be satisfied, according to our
 {\tt EKHAD}-normalization method. In this direction, the $q$-WZ pair required to prove the following Theorem is determined by 
 applying {\tt EKHAD}-normalization to a $q$-analogue of a generalization of \eqref{freelowerquarter}, namely 
\begin{equation}\label{prototypeqF}
 F(n, k) := q^{k} \left[ \begin{matrix} 
 q^{a}, q^{b} \vspace{1mm} \\ 
 q^{n+c}, q^{n+d} \end{matrix} \, \Bigg| \, q \right]_{k}, 
\end{equation}
 for free parameters $a$, $b$, $c$, and $d$. 

\begin{theorem}\label{theoremquarter}
 Letting

 \ 
 
\noindent $ p(n) = q^{a+b+c+1}-q^{a+b+c+d+n} + q^{a+b+d+1} -  q^{a+c+d+n+1}-q^{b+c+d+n+1}+q^{2 c+2 d+3 n}$, 

 \ 

\noindent the relation 
\begin{multline*}
 (q^{a+b} - q^{c+d}) (q^{a+b+1} - q^{c+d}) 
 {}_{3}\phi_{2}\!\!\left[ \begin{matrix} 
 q^{a}, q^{b}, q \vspace{1mm}\\ 
 q^{c}, q^{d} \end{matrix} \ \Bigg| \ q; q 
 \right] = \\ 
 \sum_{n=0}^{\infty} q^{n} 
 \left[ \begin{matrix} 
 q^{-a+c}, q^{-b+c}, q^{-a+d}, q^{-b+d} \vspace{1mm} \\ 
 q^{c}, q^{d} \end{matrix} \, \Bigg| \, q \right]_{n} 
 \frac{p(n)}{ 
 \big( q^{1-a-b+c+d};q^2 \big)_{n} \big( q^{2-a-b+c+d};q^2 \big)_{n} }
 \end{multline*}
 holds if $\big| \frac{q^{a + b + 1}}{q^{c + d}} \big| < 1 < |q|$. 
\end{theorem} 

\begin{proof}
 For convenience, we write $\alpha = q^a$, $\beta = q^b$, $\gamma = q^c$,  and $\delta = q^d$.  Define  $$ \overline{F}(n, k; q) :=  
 q^{k} \left[ \begin{matrix} 
 \alpha, \beta \vspace{1mm} \\ 
 \gamma q^{n}, \delta q^{n} \end{matrix} \, \Bigg| \, q \right]_{k} 
 \left[ \begin{matrix} 
 \frac{\gamma}{\alpha}, \frac{\gamma}{\beta}, \frac{\delta}{\alpha}, \frac{\delta}{\beta} \vspace{1mm} \\ 
 \gamma, \delta \end{matrix} \, \Bigg| \, q \right]_{n} \frac{1}{ \left( \frac{\gamma \delta}{\alpha \beta q}; q \right)_{2n} },  $$  i.e., as the 
 {\tt EKHAD}-normalized version of \eqref{prototypeqF}. We then let 
\begin{multline*}
 \overline{R}(n, k; q) = \\ q^{n-k} \frac{ -\alpha \beta \gamma \delta q^{k+n} + \alpha \beta \gamma 
 q + \alpha \beta \delta q - \alpha \gamma \delta q^{n+1} - \beta \gamma \delta q^{n+1} 
 + \gamma^{2} \delta^{2} q^{3n+k} }{(\alpha \beta - \gamma \delta q^{2n}) (\alpha \beta q - \gamma \delta q^{2n})}. 
\end{multline*}
 We then define 
 $\overline{G}(n, k; q) := \overline{R}(n, k; q) \overline{F}(n, k; q)$. We may thus verify that 
 the $q$-WZ difference equation 
\begin{equation}\label{firstdiff}
 \overline{F}(n+1, k; q) - \overline{F}(n, k; q) = \overline{G}(n, k+1; q) - \overline{G}(n, k; q) 
\end{equation}
 holds. The telescoping argument from Section \ref{sectionEKHAD}
 applied to \eqref{firstdiff} then 
 gives us that 
\begin{equation}\label{kbarN}
 \sum_{k=0}^{\overline{k} - 1} \left( \overline{F}(N,k;q) - \overline{F}(0,k;q) \right) 
 = \sum_{n=0}^{N-1} \left( \overline{G}(n,\overline{k};q) - \overline{G}(n,0;q) \right). 
\end{equation}
 For $|q| > 1$, we have that 
\begin{multline*}
 \lim_{N \to \infty} \frac{ \overline{F}(N+1,k;q) }{\overline{F}(N,k;q)} 
 = \\
 \lim_{N \to \infty} \frac{ (\gamma q^{N} - \alpha) (\gamma q^{N} - \beta) (\delta q^{N} - \alpha) (\delta q^{N} - 
 \beta) 
 q }{ (1 - \gamma q^{N+k}) (1 - \delta q^{N + k}) (\gamma \delta q^{2N} - \alpha \beta q) (\gamma \delta q^{2N} - 
 \alpha \beta) } = 0, 
\end{multline*} 
 so that $\lim_{N \to \infty} \overline{F}(N,k;q) = 0$. Similarly, and again for $|q| > 1$, we find that 
 $$ \lim_{n \to \infty} \frac{ \overline{G}(n+1,k;q) }{\overline{G}(n,k;q)} 
 = \lim_{n \to \infty} \frac{ \overline{R}(n+1,k;q) }{\overline{R}(n,k;q)} \lim_{n \to \infty} \frac{\overline{F}(n+1,k;q)}{\overline{F}(n,k;q)} = 0, $$
 noting that $ \lim_{n \to \infty} \frac{ \overline{R}(n+1,k;q) }{\overline{R}(n,k;q)} = 1$. 
 Letting $N \to \infty$ in \eqref{kbarN}, we obtain that 
\begin{equation}\label{bothconverge}
 \sum_{k=0}^{\overline{k} - 1} \overline{F}(0, k; q) = \sum_{n=0}^{\infty} \overline{G}(n, 0; q) - \sum_{n = 
  0}^{\infty} \overline{G}(n, \overline{k}; q). 
\end{equation}
 Since $$ \lim_{k \to \infty} \frac{ \overline{F}(0, k+1;q) }{\overline{F}(0, k;q)} = \frac{\alpha \beta q}{\gamma \delta}, $$
 we require the condition $\big| \frac{\alpha \beta q}{\gamma \delta} \big| < 1$
 to ensure the convergence of both sides of \eqref{bothconverge}, so we assume that this condition holds. 
 From the definition of $\overline{R}(n, k;q)$, we find that 
\begin{multline*}
 \left| \overline{R}(n, k; q) \right| = 
 \Bigg| \frac{1}{ \left( \alpha \beta q^{-2n} - \gamma \delta \right) 
 \left( \alpha \beta q^{-2n+1} - \gamma \delta \right)} 
 \big( -\alpha \beta \gamma \delta q^{-2n} + \\ \alpha \beta \gamma q^{-3n - k + 1} 
 + \alpha \beta \delta q^{-3n-k+1} - \alpha \gamma \delta q^{-2n - k + 1} - \beta \gamma \delta q^{-2n-k+1}
 + \gamma^{2} \delta^{2} \big) \Bigg|. 
\end{multline*}
 Consequently, we have that  $$ \left| \overline{R}(n, k; q) \right|   \leq \frac{ |\alpha \beta \gamma \delta| + |\alpha \beta \gamma q|      + 
  |\alpha \beta \delta q| + |\alpha \gamma \delta q| + |\beta \gamma \delta q|     + |\gamma^{2} \delta^{2}| }{ \gamma^{2} \delta^{2}  
  \left| 1 - \frac{\alpha \beta q}{\gamma \delta} q^{-2n-1} \right|    \left| 1 - \frac{\alpha \beta q}{\gamma \delta} q^{-2n} \right| }, $$   and  
    we can conclude that    $ \big| \overline{R}(n, k; q) \big| \leq M$,     for a value $M$ independent of $n$ and $k$,  with 
\begin{equation}\label{boundGwithF}
 \left| \overline{G}(n, k;q) \right| =    \left| \overline{R}(n, k;q) \right| \, \left| \overline{F}(n, k; q) \right|    \leq M \, \left| \overline{F}(n, k;  
  q) \right| 
\end{equation}
 for the same value $M$. Since 
\begin{equation}\label{aktoadd}
 \left| \frac{\overline{F}(n, k+1;q)}{\overline{F}(n,k;q)} \right|    = \left| q \frac{ \left( 1 - \alpha q^{k} \right) \left( 1 - \beta q^{k} 
 \right) }{ \left( 1 - \gamma q^{n+k} \right) \left( 1 - \delta q^{n+k} \right) } \right|, 
\end{equation}
 by writing   $$ \mathcal{A}_{k} = q^{-2n} \frac{\left( 1 - \alpha^{-1} q^{-k} \right) 
 \left( 1 - \beta^{-1} q^{-k} \right) }{ \left( 1 - \gamma^{-1} q^{-n-k} \right) \left( 1 - 
 \delta^{-1} q^{-n-k} \right) }, $$
 this leads us to rewrite \eqref{aktoadd} so that 
\begin{equation*}
 \left| \frac{\overline{F}(n, k+1;q)}{\overline{F}(n,k;q)} \right| 
 = \left| \frac{\alpha \beta q}{\gamma \delta} \mathcal{A}_{k} \right|, 
\end{equation*}
 recalling that $|q| > 1$.  Since $\lim_{k \to \infty} \mathcal{A}_{k} = q^{-2n}$,   we can conclude, again using the condition that $\big|  
 \frac{\alpha \beta q}{\gamma \delta} \big| < 1$,  that 
\begin{equation*}
 \left| \frac{\overline{F}(n, k+1;q)}{\overline{F}(n,k;q)} \right|  < 1. 
\end{equation*}
 Recalling \eqref{boundGwithF}, we have, for $k > \overline{k}_{0}$, that 
\begin{align*}
 \left| \overline{F}(n, k; q) \right| & \leq \rho^{k - \overline{k}_{0}} \left| \overline{F}(n, \overline{k}_{0}; q) \right| \ \text{and} \\ 
 \left| \overline{G}(n, k;q) \right| & \leq \rho^{k - \overline{k}_{0}} M \left| \overline{F}(n, \overline{k}_{0}; q) \right|, 
\end{align*}
 for a positive value $\rho$. 
 So, for $\overline{k} > \overline{k}_{0}$, we find that
\begin{equation}\label{kbarineq}
 \sum_{n=0}^{\infty} \left| \overline{G}(n, \overline{k};q) \right| 
 \leq M \, \rho^{\overline{k} - \overline{k}_{0}} \sum_{n=0}^{\infty} \left| \overline{F}(n, \overline{k}_{0}; q) \right|, 
\end{equation} 
 with the right-hand side of \eqref{kbarineq} approaching $0$
 as $\overline{k} \to \infty$. 
 So, by letting $\overline{k} \to \infty$ in \eqref{bothconverge}, 
 we obtain that 
 $$ \sum_{k=0}^{\infty} \overline{F}(0, k;q) = \sum_{n=0}^{\infty} \overline{G}(n,0;q), $$
 which is equivalent to the desired result. 
\end{proof}

\begin{example}
 Setting $(a, b, c, d) = \big( \frac{1}{2}, \frac{1}{2}, 2, 2 \big)$ in Theorem \ref{theoremquarter}, 
 we obtain the new $q$-analogue of the Ramanujan series for $\frac{1}{4}$ 
 highlighted as a motivating result in Section \ref{subRamanujan}. 
\end{example}

\begin{example}
 Setting $(a, b, c, d) = \big( \frac{1}{2}, \frac{1}{2}, \frac{3}{2}, \frac{3}{2} \big)$ 
 in Theorem \ref{theoremquarter}, we obtain 
 the first new $q$-analogue of the Fabry--Guillera formula highlighted 
 as a motivating result in Section \ref{subFabryGuillera}. 
\end{example}

\begin{example}
 Setting $(a, b, c, d) = \big( 1, 1, 2, 2 \big)$ 
 in Theorem \ref{theoremquarter}, we obtain 
 the new $q$-analogue of the Ap\'{e}ry series 
 formula highlighted as a motivating result in Section \ref{subApery}. 
\end{example}
 
\begin{example}
 Setting $(a, b, c, d) = 
 \big( \frac{1}{2}, 1, \frac{3}{2}, \frac{3}{2} \big)$ in Theorem \ref{theoremquarter}, we obtain the $q$-analogue 
\begin{multline*}
 (1-q^3) \sum_{n=0}^{\infty} q^{n} 
 \left[ \begin{matrix} 
 q^2 \vspace{1mm} \\ 
 q^3 \end{matrix} \, \Bigg| \, q^2 \right]_{n} \frac{q^{2n}}{1-q^{2n+1}} = \\ 
 \sum_{n=0}^{\infty} q^{2n} 
 \frac{(q^2;q^2)_{n}^{2}}{(q^5;q^4)_{n} (q^7;q^4)_{n}} \frac{2 - q^{2n+2} - q^{4n+3}}{1-q^{2n+1}} 
 \end{multline*}
 of the formula 
\begin{equation*}
 6 G = 
 \sum_{n = 0}^{\infty} 
 \left( \frac{1}{4} \right)^{n} \left[ \begin{matrix} 
 1, 1 
 \vspace{1mm} \\ 
 \frac{5}{4}, \frac{7}{4} 
 \end{matrix} \right]_{n} \frac{6n+5}{2n+1} 
 \end{equation*}
 proved by Chu and Zhang \cite[Example 84]{ChuZhang2014}. 
\end{example}

\begin{example}
 Setting $(a, b, c, d) = \big( -\frac{1}{2}, \frac{1}{6}, \frac{1}{3}, 1 \big)$ in Theorem \ref{theoremquarter}, we obtain a $q$-analogue 
 of the Ramanujan-type formula 
\begin{equation*}
 \frac{8\sqrt{3}}{3} = 
 \sum_{n = 0}^{\infty} 
 \left( \frac{1}{4} \right)^{n} \left[ \begin{matrix} 
 \frac{1}{6}, \frac{5}{6}, \frac{3}{2} 
 \vspace{1mm} \\ 
 \frac{1}{3}, 1, \frac{4}{3} 
 \end{matrix} \right]_{n} (18n + 1). 
 \end{equation*}
\end{example}

\begin{example}
 Setting $(a, b, c, d) = \big( \frac{1}{3}, \frac{1}{3}, 1, 1 \big)$ in Theorem \ref{theoremquarter}, we obtain a $q$-analogue of 
 the Ramanujan-type formula 
\begin{equation*}
 \frac{3 \Gamma^{3}\left( \frac{1}{3} \right) }{2 \pi^2} = 
 \sum_{n = 0}^{\infty} 
 \left( \frac{1}{4} \right)^{n} \left[ \begin{matrix} 
 \frac{2}{3}, \frac{2}{3}, \frac{2}{3} 
 \vspace{1mm} \\ 
 1, 1, \frac{7}{6} 
 \end{matrix} \right]_{n} (9 n + 2). 
 \end{equation*}
\end{example}

\begin{example}
 Setting $(a, b, c, d) = 
 \big( \frac{1}{6}, \frac{1}{6}, 1, 1 \big)$ in Theorem \ref{theoremquarter}, we obtain a $q$-analogue 
 of the Ramanujan-type formula 
\begin{equation*}
 \frac{2^{4/3} \sqrt{3} \Gamma^{3}\left( \frac{1}{3} \right)}{\pi^2} = 
 \sum_{n = 0}^{\infty} 
 \left( \frac{1}{4} \right)^{n} \left[ \begin{matrix} 
 \frac{5}{6}, \frac{5}{6}, \frac{5}{6} 
 \vspace{1mm} \\ 
 1, 1, \frac{4}{3} 
 \end{matrix} \right]_{n} (18n+5). 
 \end{equation*}
\end{example}

\begin{example}
 Setting $(a, b, c, d) = \big( \frac{1}{3}, -\frac{1}{2}, 1, \frac{1}{6} \big)$ in Theorem \ref{theoremquarter}, we obtain a $q$-analogue of 
 the Ramanujan-type formula 
\begin{equation*}
 -\frac{ \Gamma^{3}\left( \frac{1}{6} \right) }{3 \, 2^{2/3} \pi^{3/2} } = 
 \sum_{n = 0}^{\infty} 
 \left( \frac{1}{4} \right)^{n} \left[ \begin{matrix} 
 -\frac{1}{6}, \frac{2}{3}, \frac{3}{2} 
 \vspace{1mm} \\ 
 \frac{1}{6}, 1, \frac{7}{6} 
 \end{matrix} \right]_{n} (18 n - 1). 
 \end{equation*}
\end{example}

\subsection{$q$-analogues for the convergence rate $-\frac{1}{4}$}
 By analogy with the input function in \eqref{prototypeqF}, we set 
\begin{equation}\label{inputnegquart}
 F(n, k) := q^{k} \left[ \begin{matrix} 
 q^{a}, q^{b+n} \vspace{1mm} \\ 
 q^{n+c}, q^{2n+d} \end{matrix} \, \Bigg| \, q \right]_{k}, 
\end{equation}
 and apply {\tt EKHAD}-normalization to \eqref{inputnegquart}, 
 and, by mimicking the proof of 
 Theorem \ref{theoremquarter}, this provides a powerful method for generating
 $q$-analogues of accelerated series of convergence rate $-\frac{1}{4}$. 

\begin{theorem}\label{theoremnegquart}
 For the $q$-polynomial

 \ 

\noindent $p(n) = q^{n (a-b-1)} \big(-q^{a+b+c+d+3 n}-q^{a+b+c+d+3 n+1}+q^{a+b+c+2 d+5 n}-q^{a+b+2 d+4 n+1}+q^{2 a+b+1}+q^{a+c+2 d+4 n+1}-q^{2 a+d+n+1}+q^{a+2
 d+3 n+1}+q^{b+c+2 d+5 n+1}-q^{c+2 d+4 n+1}+q^{2 c+2 d+5 n}-q^{2 c+3 d+7 n}\big)$, 

 \ 

\noindent the relation 
\begin{multline*}
 (1-q^d) (q^{c+d} - q^{a+b}) (q^{c+d} - q^{a+b+1}) 
 {}_{3}\phi_{2}\!\!\left[ \begin{matrix} 
 q^{a}, q^{b}, q \vspace{1mm}\\ 
 q^{c}, q^{d} \end{matrix} \ \Bigg| \ q; q 
 \right] \\ 
 = \sum_{n=0}^{\infty} (-1)^{n} q^{-\binom{n}{2}} 
 \left[ \begin{matrix} 
 q^{-a+d}, q^{-a+d+1} \vspace{1mm} \\ 
 q^{d+1}, q^{d+2}, q^{-a-b+c+d+1}, q^{-a-b+c+d+2} \end{matrix} \, \Bigg| \, q^2 \right]_{n} \times \\ 
 \left[ \begin{matrix} 
 q^{b}, q^{-a+c}, q^{-b+d} \vspace{1mm} \\ 
 q^{c} \end{matrix} \, \Bigg| \, q \right]_{n} \, p(n)
\end{multline*}
 holds if $\big| \frac{q^{a + b + 1}}{q^{c + d}} \big| < 1 < |q|$. 
\end{theorem} 

\begin{proof}
 By setting $F(n, k)$ as in \eqref{inputnegquart},   and by, as above, applying {\tt EKHAD}-normalization to $F(n, k)$, 
 we obtain a $q$-WZ pair 
 $(\overline{F}, \overline{G})$. Omitting details, 
 by mimicking the proof of 
 Theorem \ref{theoremquarter}, 
 we arrive at the desired identity 
 $$ \sum_{k=0}^{\infty} \overline{F}(0, k;q) = \sum_{n=0}^{\infty} \overline{G}(n,0;q), $$
 subject to the given conditions, 
 and this is equivalent to the desired result. 
\end{proof}

\begin{example}
 Setting $(a, b, c, d) = \big( \frac{1}{2}, 1, \frac{3}{2}, \frac{3}{2} \big)$ in Theorem \ref{theoremnegquart}, we obtain 
 the $q$-analogue 
\begin{multline*}
 q^2 (1-q) (1-q^3)^{2} 
 \sum_{n=0}^{\infty} 
 \left[ \begin{matrix} 
 q^2 \vspace{1mm} \\ 
 q^3 \end{matrix} \, \Bigg| \, q^2 \right]_{n} 
 \frac{q^{2n}}{1 - q^{2n+1}} = 
 \sum_{n=0}^{\infty} (-1)^{n} 
 q^{-n^2-2n} \times \\ 
 \left[ \begin{matrix} 
 q^2, q^4 \vspace{1mm} \\ 
 q^5, q^5, q^7, q^7 \end{matrix} \, \Bigg| \, q^4 \right]_{n} (q^2;q^2)_{n}^{2} \, 
 \big( 1 - q^{6 n+5} - 2 q^{8 n+5}+q^{10 n+7}+q^{12 n+8} \big) 
\end{multline*}
 of the formula 
\begin{equation*}
 18 G = 
 \sum_{n = 0}^{\infty} 
 \left( -\frac{1}{4} \right)^{n} \left[ \begin{matrix} 
 \frac{1}{2}, 1, 1, 1 
 \vspace{1mm} \\ 
 \frac{5}{4}, \frac{5}{4}, \frac{7}{4}, \frac{7}{4} 
 \end{matrix} \right]_{n} \big( 40	 n^2 + 56n + 19 \big) 
 \end{equation*}
 due to Chu and Zhang \cite[Example 85]{ChuZhang2014}. 
\end{example}

\begin{example}
 Setting $(a, b, c, d) = \big( 1, 1, 2, 2 \big)$ in Theorem \ref{theoremnegquart}, we obtain 
 the $q$-analogue 
\begin{multline*}
 q (1-q^2) 
 \sum_{n=0}^{\infty} 
 \frac{q^n}{(1-q^{n+1})^2} = \\ \sum_{n=0}^{\infty} (-1)^{n} 
 q^{-\binom{n}{2}-n} 
 \frac{ (q;q)_{n} (q^2;q)_{n} }{ (q^3;q^2)_{n} (q^4;q^2)_{n} }
 \frac{1 + q^{n+1} + q^{2n+2} - 2 q^{4n+3} - q^{5n+4}}{ (1-q^{2n+1}) (1-q^{2n+2}) }
\end{multline*}
 of the formula 
\begin{equation}\label{20251q129wwww411qqqqqqqqqqqqq103AM1A}
 \frac{2\pi^2}{3} = 
 \sum_{n = 0}^{\infty} 
 \left( -\frac{1}{4} \right)^{n} \left[ \begin{matrix} 
 1 
 \vspace{1mm} \\ 
 \frac{3}{2} \end{matrix} \right]_{n} \frac{10 n + 7}{(n+1)(2n+1)}, 
\end{equation}
 proved by Chu via well poised $\Omega$-sums \cite[Example 33]{Chu2021Omega}. 
\end{example}

\begin{example}
 Setting $(a, b, c, d) = \big( 1, 1, \frac{3}{2}, 2 \big)$ in Theorem \ref{theoremnegquart}, we obtain the new $q$-analogue
\begin{multline*}
 q^2 (1-q^3) 
 \sum_{n=0}^{\infty} 
 \left[ \begin{matrix} 
 q^2 \vspace{1mm} \\ 
 q^3 \end{matrix} \, \Bigg| \, q^2 \right]_{n} 
 \frac{q^{2n}}{1 - q^{2n+2}} = \sum_{n=0}^{\infty} (-1)^{n} 
 q^{-n^2-n} 
 \frac{(q^2;q^2)_{n}^{2}}{(q^5;q^4)_{n}(q^7;q^4)_{n} } \times \\ 
 \frac{1 + q^{2 n+1}+q^{4 n+2}+q^{6 n+4}-q^{6 n+5}-q^{8 n+5}-q^{8 n+6}-q^{10 n + 
 7} }{\left(1-q^{2 n+1}\right) \left(1 + q^{2 n+1}\right) \left(1 + q^{2
 n+2}\right)}, 
\end{multline*}
 of the formula 
\begin{equation}\label{20q25q1qq12q9wrwrrwwr4r1e1e1e0e3wAwqM1qA}
 \frac{3\pi^2}{8} = 
 \sum_{n = 0}^{\infty} 
 \left( -\frac{1}{4} \right)^{n} \left[ \begin{matrix} 
 1, 1 
 \vspace{1mm} \\ 
 \frac{5}{4}, \frac{7}{4} 
 \end{matrix} \right]_{n} \frac{5n+4}{2n+1}, 
 \end{equation}
 proved by Chu via well poised $\Omega$-sums \cite[Example 32]{Chu2021Omega}. 
 Setting $(a, b, c, d) = \big( \frac{1}{2}, \frac{1}{2}, \frac{3}{2}, \frac{3}{2} \big)$ a further 
 $q$-analogue of the same Chu formula, namely 
\begin{multline*}
 (1-q)^2 q (1+q+q^2) 
 \sum_{n=0}^{\infty} 
 \frac{q^{2n}}{(1-q^{2n+1})^{2} } = 
 \sum_{n=0}^{\infty} (-1)^{n} 
 q^{-n(n+1)} 
 \frac{ \left( q^2; q^2 \right)_{n}^{2} }{ \left( q^{5};q^{4} \right)_{n} 
 \left( q^{7};q^{4} \right)_{n} } \times \\ 
 \frac{-q^{2 n+1}-q^{4 n+2}-q^{6 n+4}+q^{6 n+5}+q^{8 n+5}+q^{8 n+6}+q^{10 n+7}-1}{\left(q^{2 n+1}-1\right) \left(q^{2 n+1}+1\right) \left(q^{2
 n+2}+1\right)}. 
\end{multline*}
\end{example}
 
\begin{example}
 Setting $(a, b, c, d) = \big( \frac{1}{6}, \frac{1}{6}, \frac{2}{3}, 1 \big)$ in Theorem \ref{theoremnegquart}, for the $q$-polynomial

 \ 

\noindent $ p(n) = \big(1-q^{2 n+1}\big) 
 \big(1 + q^{2 n+1}+q^{4 n+2}\big) 
 \big(1 + q^{6 n+1}+q^{6 n+3}-q^{6 n+5}+q^{12 n+2}+q^{12 n+4}-q^{12 n+8}+q^{18
 n+5}-q^{18 n+9}+q^{18 n+10}-q^{18 n+11}-q^{24 n+10}-q^{24 n+12}-q^{30 n+13}\big), $

 \ 

\noindent we obtain the $q$-analogue 
\begin{multline*}
 q ( 1 + q) \left(1-q^6\right) \left(1-q^8\right) 
 \sum_{n=0}^{\infty} q^{6n} 
 \left[ \begin{matrix} 
 q, q \vspace{1mm} \\ 
 q^4, q^6 \end{matrix} \, \Bigg| \, q^6 \right]_{n} 
 = \\ \sum_{n=0}^{\infty} (-1)^{n} q^{-3n^2-3n} 
 \left[ \begin{matrix} 
 q^5, q^{11} \vspace{1mm} \\ 
 q^{12}, q^{14}, q^{18}, q^{20} \end{matrix} \, \Bigg| \, q^{12} \right]_{n} 
 \left[ \begin{matrix} 
 q^{3}, q^{5}, q^{7} \vspace{1mm} \\ 
 q^{4} \end{matrix} \, \Bigg| \, q^6 \right]_{n} \, p(n)
\end{multline*}
 of the Ramanujan-type formula 
\begin{equation*}
 \frac{64\sqrt{3}}{3} = 
 \sum_{n = 0}^{\infty} 
 \left( -\frac{1}{4} \right)^{n} \left[ \begin{matrix} 
 \frac{5}{12}, \frac{5}{6}, \frac{11}{12} 
 \vspace{1mm} \\ 
 \frac{2}{3}, 1, \frac{5}{3} 
 \end{matrix} \right]_{n} (60 n + 43) 
 \end{equation*}
 introduced and proved by Chu in 2021 \cite{Chu2021Omega}
 via well poised $\Omega$-sums. 
\end{example}

\begin{example}
 Setting $(a, b, c, d) = \big( \frac{1}{3}, \frac{1}{6}, 1, \frac{5}{6} \big)$ in Theorem \ref{theoremnegquart}, we obtain 
 a $q$-analogue of the Ramanujan-type formula 
\begin{equation*}
 \frac{20 \sqrt[3]{2} }{3} = 
 \sum_{n = 0}^{\infty} 
 \left( -\frac{1}{4} \right)^{n} \left[ \begin{matrix} 
 \frac{1}{4}, \frac{2}{3}, \frac{3}{4} 
 \vspace{1mm} \\ 
 \frac{11}{12}, 1, \frac{17}{12}
 \end{matrix} \right]_{n} (20 n + 9) 
 \end{equation*}
 given previously via a different method
 in our past work \cite{Campbellunpublished}. 
\end{example}

\begin{example}
 Setting $(a, b, c, d) = \big( \frac{1}{2}, 1, \frac{3}{2}, 2 \big)$ in Theorem \ref{theoremnegquart}, we obtain a $q$-analogue of the 
 Ramanujan-type formula 
\begin{equation}\label{20251q129411ee1e0e3eqAqMq1qA}
 16 \log 2 = 
 \sum_{n = 0}^{\infty} 
 \left( -\frac{1}{4} \right)^{n} \left[ \begin{matrix} 
 \frac{3}{4}, 1, \frac{5}{4} 
 \vspace{1mm} \\ 
 \frac{3}{2}, \frac{3}{2}, \frac{3}{2}
 \end{matrix} \right]_{n} (20 n + 13) 
 \end{equation}
 proved via a different method in our past work \cite{Campbellunpublished}. 
\end{example}

 \begin{example}
 Setting $(a, b, c, d) = \big( \frac{3}{4}, \frac{1}{4}, \frac{3}{2}, 1 \big)$ in Theorem \ref{theoremnegquart}, we obtain 
 a $q$-analogue of the Ramanujan-type formula 
\begin{equation*}
 16 \sqrt{2} = 
 \sum_{n = 0}^{\infty} 
 \left( -\frac{1}{4} \right)^{n} \left[ \begin{matrix} 
 \frac{1}{8}, \frac{5}{8}, \frac{3}{4} 
 \vspace{1mm} \\ 
 1, \frac{3}{2}, \frac{3}{2} 
 \end{matrix} \right]_{n} (40 n + 23) 
 \end{equation*}
 proved via a different method in our past work \cite{Campbellunpublished}. 
\end{example}

 \begin{example}
 Setting $(a, b, c, d) = \big( \frac{3}{4}, 1, \frac{3}{2}, \frac{3}{2} \big)$ in Theorem \ref{theoremnegquart}, we obtain 
 a $q$-analogue of the Ramanujan-type formula 
\begin{equation*}
 \frac{5 \Gamma^{4}\left( \frac{1}{4} \right) }{16 \pi} = 
 \sum_{n = 0}^{\infty} 
 \left( -\frac{1}{4} \right)^{n} \left[ \begin{matrix} 
 \frac{3}{8}, \frac{7}{8}, 1 
 \vspace{1mm} \\ 
 \frac{9}{8}, \frac{5}{4}, \frac{13}{8} 
 \end{matrix} \right]_{n} (40 n + 19). 
 \end{equation*}
\end{example}

\begin{example}
 Setting $(a, b, c, d) = \big( \frac{1}{4}, \frac{1}{4}, 1, 1 \big)$ in Theorem \ref{theoremnegquart}, we obtain 
 a $q$-analogue of the Ramanujan-type formula 
\begin{equation*}
 \frac{8 \Gamma^{2}\left( \frac{1}{4} \right) }{\pi^{3/2}} = 
 \sum_{n = 0}^{\infty} 
 \left( -\frac{1}{4} \right)^{n} \left[ \begin{matrix} 
 \frac{3}{8}, \frac{3}{4}, \frac{7}{8} 
 \vspace{1mm} \\ 
 1, 1, \frac{3}{2}
 \end{matrix} \right]_{n} (40 n + 21). 
 \end{equation*}
\end{example}

 \begin{example}
 Setting $(a, b, c, d) = \big( \frac{1}{2}, \frac{5}{6}, 1, 2 \big)$ in Theorem \ref{theoremnegquart}, we obtain 
 a $q$-analogue of the Ramanujan-type formula 
\begin{equation*}
 \frac{2^{\frac{25}{3}} \pi}{3 \Gamma^{3}\left( \frac{1}{3} \right) } = 
 \sum_{n = 0}^{\infty} 
 \left( -\frac{1}{4} \right)^{n} \left[ \begin{matrix} 
 \frac{3}{4}, \frac{7}{6}, \frac{5}{4} 
 \vspace{1mm} \\ 
 1, \frac{4}{3}, 2 
 \end{matrix} \right]_{n} (20 n + 21). 
 \end{equation*}
\end{example}

\begin{example}
 Setting $(a, b, c, d) = \big( \frac{3}{4}, \frac{1}{2}, 1, \frac{3}{2} \big)$ in Theorem \ref{theoremnegquart}, we obtain 
 a $q$-analogue of the Ramanujan-type formula 
\begin{equation*}
 \frac{15 \Gamma^{2}\left( \frac{1}{4} \right) }{28 \sqrt{2\pi}} = 
 \sum_{n = 0}^{\infty} 
 \left( -\frac{1}{4} \right)^{n} \left[ \begin{matrix} 
 \frac{3}{8}, \frac{1}{2}, \frac{15}{8} 
 \vspace{1mm} \\ 
 \frac{9}{8}, \frac{13}{8}, \frac{7}{4} 
 \end{matrix} \right]_{n} (5 n + 3). 
 \end{equation*}
\end{example}

 \begin{example}
 Setting $(a, b, c, d) = \big( \frac{3}{2}, \frac{1}{6}, 2, 1 \big)$ in Theorem \ref{theoremnegquart}, we obtain 
 a $q$-analogue of the Ramanujan-type formula 
\begin{equation*}
 \frac{ 2^{\frac{14}{3}} \sqrt{3} \Gamma^{3}\left( \frac{1}{3} \right) }{5 \pi^2} = 
 \sum_{n = 0}^{\infty} 
 \left( -\frac{1}{4} \right)^{n} \left[ \begin{matrix} 
 -\frac{1}{4}, \frac{1}{4}, \frac{5}{6} 
 \vspace{1mm} \\ 
 1, \frac{5}{3}, 2 
 \end{matrix} \right]_{n} (20 n + 17). 
 \end{equation*}
\end{example}

\begin{example}
 Setting $(a, b, c, d) = \big( -\frac{1}{2}, \frac{1}{4}, 1, \frac{1}{2} \big)$ in Theorem \ref{theoremnegquart}, we obtain 
 a $q$-analogue of the Ramanujan-type formula 
\begin{equation*}
 \frac{21 \pi^{3/2} }{2 \sqrt{2} \Gamma^{2}\left( \frac{1}{4} \right)} = 
 \sum_{n = 0}^{\infty} 
 \left( -\frac{1}{4} \right)^{n} \left[ \begin{matrix} 
 \frac{1}{2}, \frac{5}{4}, \frac{3}{2} 
 \vspace{1mm} \\ 
 \frac{3}{4}, \frac{11}{8}, \frac{15}{8} 
 \end{matrix} \right]_{n} (5 n + 4). 
 \end{equation*}
\end{example}

\begin{example}
 Setting $(a, b, c, d) = \big( \frac{1}{2}, \frac{1}{6}, 1, 1 \big)$ in Theorem \ref{theoremnegquart}, we obtain 
 a $q$-analogue of the Ramanujan-type formula 
\begin{equation*}
 \frac{2^{11/3} \Gamma^{3}\left( \frac{1}{3} \right)}{ 5 \sqrt{3} \pi^2 } = 
 \sum_{n = 0}^{\infty} 
 \left( -\frac{1}{4} \right)^{n} \left[ \begin{matrix} 
 \frac{1}{4}, \frac{3}{4}, \frac{5}{6} 
 \vspace{1mm} \\ 
 1, 1, \frac{5}{3}
 \end{matrix} \right]_{n} (4 n + 3). 
 \end{equation*}
\end{example}

\subsection{Further $q$-analogues for the convergence rate $\frac{1}{4}$}\label{suqqqqbqqquqqaqrqtqer}
 For brevity, we henceforward omit full proofs of multiparameter ${}_{3}\phi_{2}$-series obtained
 by {\tt EKHAD}-normalization, 
 leaving it to the reader to adapt 
 the above proof of 
 Theorem \ref{theoremquarter}. 
 By applying {\tt EKHAD}-normalization
 starting with 
\begin{equation}\label{Fsecondquarter}
 F(n, k) := q^{k} \left[ \begin{matrix} 
 q^{n + a}, q^{n+b} \vspace{1mm} \\ 
 q^{n+c}, q^{2n+d} \end{matrix} \, \Bigg| \, q \right]_{k}, 
\end{equation}
 and by mimicking the proof of 
 Theorem \ref{theoremquarter}, 
 the resulting tranformation of a ${}_{3}\phi_{2}$-series of the form 
 $$ {}_{3}\phi_{2}\!\!\left[ \begin{matrix} 
 q^{a}, q^{b}, q \vspace{1mm}\\ 
 q^{c}, q^{d} \end{matrix} \ \Bigg| \ q; q 
 \right], $$
 by analogy with Theorems \ref{theoremquarter} 
 and \ref{theoremnegquart} and resulting from a relation of the form 
 $$ \sum_{k=0}^{\infty} \overline{F}(0, k;q) = \sum_{n=0}^{\infty} \overline{G}(n,0;q) $$
 for the $q$-WZ pair $(\overline{F}, \overline{G})$
 obtained from the application of {\tt EKHAD}-normalization to 
 \eqref{Fsecondquarter}, this gives a powerful way of obtaining 
 $q$-analogues of accelerated series of convergence rate $\frac{1}{4}$, 
 as demonstrated below. 

\begin{example}
 Setting $(a, b, c, d) = \big( 1, 1, \frac{3}{2}, 2 \big)$ after applying {\tt EKHAD}-normali-zation to \eqref{Fsecondquarter}, 
 we obtain the second new $q$-analogue of the Fabry--Guillera formula highlighted 
 as a motivating result in Section \ref{subFabryGuillera}. 
\end{example} 

\begin{example}
 Setting $(a, b, c, d) = \big( \frac{1}{4}, \frac{1}{4}, 1, \frac{3}{4} \big)$ 
 after applying {\tt EKHAD}-normali-zation to \eqref{Fsecondquarter}, 
 we obtain a $q$-analogue of the Ramanujan-type formula 
 $$ \frac{3\sqrt{2}}{4} = \sum_{n = 0}^{\infty} 
 \left( \frac{1}{4} \right)^{n} \left[ \begin{matrix} 
 \frac{1}{4}, \frac{1}{2}, \frac{1}{2} \vspace{1mm} \\ 
 \frac{7}{8}, 1, \frac{11}{8} \end{matrix} \right]_{n} (3n+1). $$
\end{example} 

\begin{example}
 Setting $(a, b, c, d) = \big( -\frac{1}{2}, \frac{1}{3}, 1, \frac{1}{6} \big)$ after applying {\tt EKHAD}-normali-zation to \eqref{Fsecondquarter}, 
 we obtain a $q$-analogue of the Ramanujan-type formula 
 $$ \frac{8 \sqrt{3}}{3} = \sum_{n = 0}^{\infty} 
 \left( \frac{1}{4} \right)^{n} \left[ \begin{matrix} 
 -\frac{1}{2}, \frac{1}{6}, \frac{5}{6} \vspace{1mm} \\ 
 \frac{2}{3}, 1, \frac{5}{3} \end{matrix} \right]_{n} (18 n + 5). $$
\end{example} 

\begin{example}
 Setting $(a, b, c, d) = \big( \frac{1}{6}, \frac{1}{2}, 1, \frac{5}{6} \big)$ after applying {\tt EKHAD}-normali-zation to \eqref{Fsecondquarter}, 
 we obtain a $q$-analogue of the Ramanujan-type formula 
 $$ \frac{5 \sqrt{3}}{2} = \sum_{n = 0}^{\infty} 
 \left( \frac{1}{4} \right)^{n} \left[ \begin{matrix} 
 \frac{1}{3}, \frac{1}{2}, \frac{2}{3} \vspace{1mm} \\ 
 \frac{11}{12}, 1, \frac{17}{12} \end{matrix} \right]_{n} (9 n + 4). $$
\end{example} 

\begin{example}
 Setting $(a, b, c, d) = \big( \frac{1}{3}, \frac{5}{6}, 1, \frac{3}{2} \big)$ after applying {\tt EKHAD}-normali-zation to \eqref{Fsecondquarter}, 
 we obtain a $q$-analogue of the Ramanujan-type formula 
 $$ \frac{27 \sqrt[3]{2} }{4} = \sum_{n = 0}^{\infty} 
 \left( \frac{1}{4} \right)^{n} \left[ \begin{matrix} 
 \frac{2}{3}, \frac{5}{6}, \frac{7}{6} \vspace{1mm} \\ 
 1, \frac{5}{4}, \frac{7}{4} \end{matrix} \right]_{n} (9 n + 7). $$
\end{example} 

\begin{example}
 Setting $(a, b, c, d) = \big( \frac{1}{3}, \frac{5}{6}, \frac{3}{2}, 1 \big)$ 
 after applying {\tt EKHAD}-normali-zation to \eqref{Fsecondquarter}, 
 we obtain a $q$-analogue of the Ramanujan-type formula 
 $$ 9 \sqrt[3]{2} = \sum_{n = 0}^{\infty} 
 \left( \frac{1}{4} \right)^{n} \left[ \begin{matrix} 
 \frac{1}{6}, \frac{2}{3}, \frac{5}{6} \vspace{1mm} \\ 
 1, \frac{3}{2}, \frac{3}{2} \end{matrix} \right]_{n} (18 n + 11). $$
\end{example} 

\begin{example}
 Setting $(a, b, c, d) = \big( \frac{1}{6}, \frac{1}{3}, 1, \frac{5}{6} \big)$ after applying {\tt EKHAD}-normali-zation to \eqref{Fsecondquarter}, 
 we obtain a $q$-analogue of the Ramanujan-type formula 
 $$ \frac{5 \sqrt[3]{2} }{6} = \sum_{n = 0}^{\infty} 
 \left( \frac{1}{4} \right)^{n} \left[ \begin{matrix} 
 \frac{1}{6}, \frac{1}{2}, \frac{2}{3} \vspace{1mm} \\ 
 \frac{11}{12}, 1, \frac{17}{12} \end{matrix} \right]_{n} (3 n + 1). $$
\end{example} 

\begin{example}
 Setting $(a, b, c, d) = \big( \frac{1}{2}, -\frac{1}{2}, \frac{3}{4}, 1 \big)$ 
 after applying {\tt EKHAD}-normali-zation to \eqref{Fsecondquarter}, 
 we obtain a $q$-analogue of the Ramanujan-type formula 
 $$ \frac{12\pi^2}{\Gamma^{4}\left( \frac{1}{4} \right)} = \sum_{n = 0}^{\infty} 
 \left( \frac{1}{4} \right)^{n} \left[ \begin{matrix} 
 -\frac{1}{2}, \frac{1}{2}, \frac{3}{2} \vspace{1mm} \\ 
 \frac{3}{4}, 1, \frac{7}{4} \end{matrix} \right]_{n} (3n+1). $$
\end{example} 

\begin{example}
 Setting $(a, b, c, d) = \big( \frac{1}{2}, -\frac{1}{2}, \frac{5}{6}, 1 \big)$ 
 after applying {\tt EKHAD}-normali-zation to \eqref{Fsecondquarter}, 
 we obtain a $q$-analogue of the Ramanujan-type formula 
 $$ \frac{ 5 \, 2^{\frac{11}{3}} \pi^3 }{ \Gamma^{6}\left( \frac{1}{3} 
 \right) } = \sum_{n = 0}^{\infty} 
 \left( \frac{1}{4} \right)^{n} \left[ \begin{matrix} 
 -\frac{1}{2}, \frac{1}{2}, \frac{3}{2} \vspace{1mm} \\ 
 \frac{5}{6}, 1, \frac{11}{6} \end{matrix} \right]_{n} (18n+7). $$
\end{example} 

\begin{example}
 Setting $(a, b, c, d) = \big( \frac{1}{2}, \frac{5}{6}, 2, 1 \big)$ 
 after applying {\tt EKHAD}-normali-zation to \eqref{Fsecondquarter}, 
 we obtain a $q$-analogue of the Ramanujan-type formula 
 $$ \frac{ 2^{\frac{19}{3}} \pi }{ \Gamma^{3}\left( \frac{1}{3} 
 \right) } = \sum_{n = 0}^{\infty} 
 \left( \frac{1}{4} \right)^{n} \left[ \begin{matrix} 
 \frac{1}{6}, \frac{1}{2}, \frac{5}{6} \vspace{1mm} \\ 
 1, \frac{5}{3}, 2 \end{matrix} \right]_{n} (18n+13). $$
\end{example} 

\begin{example}
 Setting $(a, b, c, d) = \big( -\frac{1}{2}, \frac{1}{2}, \frac{1}{3}, 1 \big)$ 
 after applying {\tt EKHAD}-normali-zation to \eqref{Fsecondquarter}, 
 we obtain a $q$-analogue of the Ramanujan-type formula 
 $$ -\frac{\Gamma^{6}\left( \frac{1}{3} \right)}{2^{\frac{5}{3}} \pi^3} = \sum_{n = 0}^{\infty} 
 \left( \frac{1}{4} \right)^{n} \left[ \begin{matrix} 
 -\frac{1}{2}, \frac{1}{2}, \frac{3}{2} \vspace{1mm} \\ 
 \frac{1}{3}, 1, \frac{4}{3} \end{matrix} \right]_{n} (18n + 1). $$
\end{example} 

\begin{example}
 Setting $(a, b, c, d) = \big( -\frac{1}{2}, \frac{1}{3}, 1, \frac{1}{6} \big)$ 
 after applying {\tt EKHAD}-normali-zation to \eqref{Fsecondquarter}, 
 we obtain a $q$-analogue of the Ramanujan-type formula 
 $$ -\frac{ \sqrt{3} \Gamma^{6}\left( \frac{1}{3} \right) }{2^{\frac{14}{3}} \pi^3} = 
 \sum_{n = 0}^{\infty} 
 \left( \frac{1}{4} \right)^{n} \left[ \begin{matrix} 
 -\frac{1}{2}, -\frac{1}{6}, \frac{2}{3} \vspace{1mm} \\ 
 \frac{7}{12}, 1, \frac{13}{12} \end{matrix} \right]_{n} (9n - 1). $$
\end{example} 

\subsection{$q$-analogues for the convergence rate $\frac{1}{64}$}
 Applying {\tt EKHAD}-normalization to 
\begin{equation}\label{Ffor64}
 F(n, k) := q^{k} \left[ \begin{matrix} 
 q^{n+a}, q^{n+b} \vspace{1mm} \\ 
 q^{2n+c}, q^{2n+d} \end{matrix} \, \Bigg| \, q \right]_{k}, 
\end{equation}
 by analogy with the preceding subsection, 
 this gives us 
 a way of obtaining $q$-analogues of 
 accelerated series of convergence rate $\frac{1}{64}$, as below. 

\begin{example}
 Setting $(a,b,c,d) = \big( 1, 1, 2, 2 \big)$ after applying {\tt EKHAD}-normali-zation to \eqref{Ffor64}, 
 this provides the $q$-analogue 
 highlighted in Section \ref{sectionZeilberger}
 of the Zeilberger formula in \eqref{displayZeilberger}. 
\end{example}

\begin{example}
 Setting $(a,b,c,d) = \big( \frac{1}{2}, \frac{1}{2}, 1, \frac{3}{2} \big)$
 after applying {\tt EKHAD}-normali-zation to 
 \eqref{Ffor64}, 
 and setting 

 \ 

\noindent $ r(n) = \big( 3 q^{2 n+2}+q^{4 n+1}+q^{4 n+3}+q^{6 n+2}+q^{6 n+4}-2 q^{8 n+4}-2 q^{8 n+6}-2 q^{10 n+5}+q^{12 n+7}+q^{14 n+8}-q-2
 \big)/\big( q^{2 n+1}+1 \big), $ 

 \ 

\noindent we obtain the $q$-analogue 
\begin{multline*}
 -(1-q)^2 (1-q^3)^{2} \sum_{k=0}^{\infty} \frac{q^{2k}}{1-q^{2k+1}} 
 \left[ \begin{matrix} 
 q \vspace{1mm} \\ 
 q^2 \end{matrix} \, \Bigg| \, q^2 \right]_{k} = \\ 
 \sum_{n=0}^{\infty} q^{2n} \frac{ \left( q;q^{2} \right)_{n}^{4} 
 \left( q^{2};q^{2} \right)_{n}^{2} }{ \left( q^{2};q^{4} \right)_{n} 
 \left( q^{4};q^{4} \right)_{n} 
 \left( q^{5};q^{4} \right)_{n}^{2} \left( q^{7};q^{4} \right)_{n}^{2} } r(n)
\end{multline*} 
 of the Chu formula 
\begin{equation*}
 \frac{9\pi}{4} = 
 \sum_{n = 0}^{\infty} 
 \left( \frac{1}{64} \right)^{n} \left[ \begin{matrix} 
 1, \frac{1}{2}, \frac{1}{2}, \frac{1}{2} 
 \vspace{1mm} \\ 
 \frac{5}{4}, \frac{5}{4}, \frac{7}{4}, \frac{7}{4}
 \end{matrix} \right]_{n} \left( 42n^3 + 75n^2+42n + 7 \right) 
 \end{equation*} 
 introduced via well-poised $\Omega$-sums \cite{Chu2021Omega}. 
\end{example}

\begin{example}
 Setting $(a,b,c,d) = \big( \frac{1}{2}, \frac{1}{6}, 1, 1 \big)$
 after applying {\tt EKHAD}-normali-zation to 
 \eqref{Ffor64}, 
 we obtain a $q$-analogue of the Ramanujan-type formula 
\begin{equation*}
 \frac{ 2^{\frac{14}{3}} \sqrt{3} \Gamma^{3}\left( \frac{1}{3} \right) }{\pi^2} = 
 \sum_{n = 0}^{\infty} 
 \left( \frac{1}{64} \right)^{n} \left[ \begin{matrix} 
 \frac{1}{2}, \frac{5}{6}, \frac{5}{6} 
 \vspace{1mm} \\ 
 1, 1, \frac{5}{3} 
 \end{matrix} \right]_{n} \left( 126 n + 85 \right). 
 \end{equation*} 
\end{example}

\subsection{$q$-analogues for the convergence rate $-\frac{1}{27}$}
 By applying {\tt EKHAD}-normali-zation to 
 \begin{equation}\label{inputfor27}
 q^{k} \left[ \begin{matrix} 
 q^{a}, q^{n+b} \vspace{1mm} \\ 
 q^{2n+c}, q^{2n+d} \end{matrix} \, \Bigg| \, q \right]_{k}, 
\end{equation}
 by analogy with the preceding section, this gives us a way of generating 
 $q$-analogues of accelerated series of convergence rate $-\frac{1}{27}$, as below. 

\begin{example}
 Setting $(a, b, c, d) = \big( \frac{1}{2}, \frac{1}{2}, \frac{3}{2}, 1 \big)$ 
 after applying {\tt EKHAD}-normali-zation to \eqref{inputfor27}, 
 and by writing

 \ 

\noindent $p(n) = 1-q^{2 n+2}+q^{4 n+1}+q^{4 n+2}+q^{4 n+3}+q^{6 n+2}-q^{6 n+5}+q^{8 n+3}-q^{10 n+6}-q^{10 n+7}-q^{12 n+7}-q^{14 n+8}$, 

 \ 

\noindent we obtain the new $q$-analogue
\begin{multline*} 
 q(1-q^5) \sum_{k=0}^{\infty} 
 q^{2k} \left[ \begin{matrix} 
 q, q \vspace{1mm} \\ 
 q^2, q^3 \end{matrix} \, \Bigg| \, q^2 \right]_{k} = \\ 
 \sum_{n=0}^{\infty} (-1)^n q^{-n-n^2} 
 \frac{ (q;q^2)_{n}^{2} (q^2;q^2)_{n} }{ (q^3;q^6)_{n} (q^7;q^6)_{n} (q^{11};q^6)_{n} } 
 \frac{p(n)}{q^{2 n+1}+q^{4
 n+2}+1} 
\end{multline*}
 of the formula 
 $$ \frac{15 \pi}{8} = \sum_{k=0}^{\infty} 
 \left( -\frac{1}{27} \right)^{k} \left[ \begin{matrix} 
 1, \frac{1}{2} \vspace{1mm} \\ 
 \frac{7}{6}, \frac{11}{6} \end{matrix} \right]_{k} (7k+6) $$
 due to Chu and Zhang \cite[Example 27]{ChuZhang2014}. 
\end{example}

\begin{example}
 Setting $(a, b, c, d) = \big( 1, 1, 2, 2 \big)$ after applying {\tt EKHAD}-normali-zation to \eqref{inputfor27}, 
 and writing 
 
 \ 

\noindent $ r(n) = (-q^{2 n+1}-q^{2 n+2}-3 q^{3 n+2}-q^{4 n+3}+q^{5 n+3}+2 q^{6 n+4}+q^{6 n+5}+2 q^{7 n+5}+q^{8 n+6} - 
 1) / (\left(q^{n+1}+1\right)
 \left(q^{n+1}+q^{2 n+2}+1\right)), $ 

 \ 

\noindent we obtain the $q$-analogue
\begin{multline*}
 -q(1-q)^{3} (1+q) \sum_{k=0}^{\infty} \frac{q^{k}}{(1-q^{k+1})^{2}} = \\ 
 \sum_{n=0}^{\infty} (-1)^{n} q^{-\frac{n(n+1)}{2}}
 \frac{ \left( q;q \right)_{n}^{3} \left( q;q^2 \right)_{n} }{ 
 \left( q^{3};q^{2} \right)_{n} \left( q^{3};q^{3} \right)_{n} \left( q^{4};q^{3} \right)_{n} 
 \left( q^{5};q^{3} \right)_{n} } r(n)
\end{multline*}
 of the formula 
 $$ \frac{\pi^2}{2} = \sum_{k=0}^{\infty} 
 \left( -\frac{1}{27} \right)^{k} \left[ \begin{matrix} 
 1, 1 \vspace{1mm} \\ 
 \frac{4}{3}, \frac{5}{3} \end{matrix} \right]_{k} \frac{7k+5}{2k+1} $$
 due to Chu and Zhang \cite[Example 25]{ChuZhang2014}. 
\end{example}

\begin{example}
 Setting $(a, b, c, d) = \big( 1, \frac{1}{2}, \frac{3}{2}, \frac{3}{2} \big)$ 
 after applying {\tt EKHAD}-normali-zation to \eqref{inputfor27}, 
 and writing 
 
 \ 

\noindent $ p(n) = q^{2 n+1}-2 q^{2 n+2}+q^{4 n+1}+q^{4 n+2}-q^{4 n+3}+q^{4 n+4}+q^{6 n+2}+2 q^{6 n+3}-q^{6 n+4}-q^{6 n+5}+q^{8 n+3}-q^{8 n+4}-2 q^{8 n+5}-q^{8
 n+6}+q^{8 n+7}-q^{10 n+5}-q^{10 n+6}-q^{10 n+7}+q^{10 n+8}-q^{12 n+6}-q^{12 n+7}-q^{12 n+8}+q^{12 n+9}-q^{14 n+8}+q^{14 n+9}+q^{14
 n+10}+q^{16 n+10}+q^{18 n+11}+1, $ 

 \ 

\noindent we obtain the $q$-analogue
\begin{multline*}
 q(1-q^5) \sum_{k=0}^{\infty} 
 \frac{q^{2k}}{1-q^{2k+1}} \left[ \begin{matrix} 
 q^{2} \vspace{1mm} \\ 
 q^{3} \end{matrix} \, \Bigg| \, q^2 \right]_{k} = 
 \sum_{n=0}^{\infty} (-1)^{n} q^{-n^2} 
 \times \\ \frac{ \left( q;q^2 \right)_{n} \left( q^2;q^2 \right)_{n}^{2} }{ 
 \left( q^{3};q^{6} \right)_{n} 
 \left( q^{7}; q^{6} \right)_{n} \left( q^{11};q^{6} \right)_{n} } 
 \frac{p(n)}{ \left( 1-q^{4n+1} \right) \left( 1 - q^{4n+3} \right) 
 \left( 1 + q^{2n+1} + q^{4n+2} \right) }
\end{multline*}
 of the formula 
 $$ 30 G = \sum_{k=0}^{\infty} 
 \left( -\frac{1}{27} \right)^{k} \left[ \begin{matrix} 
 1, 1 \vspace{1mm} \\ 
 \frac{7}{6}, \frac{11}{6} \end{matrix} \right]_{k} \frac{112k^2+192k+83}{(4k+1)(4k+3)} $$
 due to Chu and Zhang \cite[Example 29]{ChuZhang2014}. 
\end{example}

\section{Conclusion}
 Our multiparameter {\tt EKHAD}-normalization method may be applied much more broadly, using $q$-analogues of hypergeometric 
 input functions as in our past work \cite{Campbell2025Hypergeometric}. For brevity, and in view of the many new $q$-analogues of 
 previously published results from many authors highlighted 
 in Section \ref{sectionMotivating}, we leave this, along with the research areas below, to separate projects. 

 We encourage the exploration of alternatives to and variants of our $q$-WZ-based techniques given above. For example, the
 $q$-WZ-based technique employed in our past work \cite{Campbell2025further} to produce the final $q$-analogue of the
 Fabry--Guillera formula listed in Section \ref{subFabryGuillera} is not equivalent to the $q$-WZ-based techniques given in our current 
 paper, noting that this previous $q$-WZ-based technique produces series with $q$-Pochhammer symbols indexed by $\infty$
 within the summands. For example, by adapting the approach from our past research on a $q$-version of the Fabry--Guillera formula 
 \cite{Campbell2025further} and by using a recent terminating $q$-binomial identity due to 
 Berkovich in 2024 \cite{Berkovich2024}, it can be shown that 
\begin{multline*}
 \sum_{n=0}^{\infty} \frac{q^{n(n+1)}}{3^n} 
 \left( \frac{1}{ \left( -3q;q^2 \right)_{\infty} \left( -\frac{q}{3}; q^2 \right)_{\infty} } 
 - \frac{ \left( q^{-2n+1};q^2 \right)_{\infty} \left( q^{2n+3};q^2 \right)_{\infty} }{ \left( q^2;q^2 \right)_{\infty} } 
 \right) = \\ 
 - 6 q \sum_{n=0}^{\infty} \frac{q^{2n}}{(q^{2n+1} + 3) (3 q^{2n+1} +1) } \frac{ 
 \left( q^{2n+1};q^2 \right)_{\infty} \left( q^{2n+3};q^2 \right)_{\infty} }{ \left( q^{4n+2};q^2 \right)_{\infty} 
 \left( -3q;q^2 \right)_{n} \left( -\frac{q}{3};q^2 \right)_{n}}, 
\end{multline*}
 and this gives us a $q$-analogue of the $\pi$-formula
 $\frac{2 \pi}{\sqrt{3}} = \sum_{n=1}^{\infty} \frac{3^n}{n \binom{2n}{n}}$. 
 
 How could {\tt EKHAD}-normalization be applied in the context of combinatorial aspects of integer partitions and associated $q$-series? 
 Give a $q$-series identity associated with a combinatorial proof 
 of an identity on integer partitions, how could an ``accelerated'' 
 version of this $q$-series identity be obtained through {\tt EKHAD}-normalization
 or an analogue of this technique? 

\subsection*{Acknowledgements}
 The author is very thankful to Yan-Ping Mu for much help formulating the proof of Theorem \ref{theoremquarter}, 
 and thanks Kam Cheong Au and Michael Schlosser for useful feedback. 

\begin{thebibliography}{99}

\bibitem{AdamchikWagon1996}%%%%%%%%%%%%%%%q
 \textsc{V.\ Adamchik and S.\ Wagon}, 
 $\pi$: a 2000-year-old search changes direction, 
 \emph{Math.\ Edu.\ Res.} {\bf 5} (1996), 11--19. 

\bibitem{Adams1931}%%%%%%%%%%%%%%%q
 \textsc{C.\ R.\ Adams}, 
 Linear {{\(q\)}}-difference equations, 
 \emph{Bull.\ Am.\ Math.\ Soc.} {\bf 37} (1931), 361--400. 

\bibitem{Andrews1968}%%%%%%%%%%%%%%%q
 \textsc{G.\ E.\ Andrews}, 
 On {$q$}-difference equations for certain well-poised basic hypergeometric series, 
 \emph{Quart.\ J.\ Math.\ Oxford Ser.\ (2)} {\bf 19} (1968), 433--447. 

\bibitem{Apery1979}%%%%%%%%%%%%%%%q
 \textsc{R.\ Ap\'ery}, 
 Irrationalit\'e{} de {$\zeta 2$} et {$\zeta 3$}, 
 \emph{Ast\'erisque} {\bf (61)} (1979), 11--13. 

\bibitem{ArndtHaenel2001}%%%%%%%%%%%%%%%q
 \textsc{J.\ Arndt and C.\ Haenel}, 
 Pi---unleashed, 
 Springer-Verlag, Berlin (2001). 

\bibitem{Auunpublished}%%%%%%%%%%%%%%%
 \textsc{K.\ C.\ Au}, 
 $q$-analogues of Wilf--Zeilberger seeds and Ramanujan $1/\pi^k$-formulas, 
 arXiv:2403.04555 (2026). 

\bibitem{Au2025}%%%%%%%%%%%%%%%
 \textsc{K.\ C.\ Au}, 
 Wilf-{Z}eilberger seeds and non-trivial hypergeometric identities, 
 \emph{J.\ Symbolic Comput.} {\bf 130} (2025), Paper No. 102421, 34. 

\bibitem{Bailey1964}%%%%%%%%%%%%%%%
 \textsc{W.\ N.\ Bailey}, 
 Generalized hypergeometric series, 
 Stechert-Hafner, Inc., New York (1964). 

\bibitem{BaruahBerndtChan2009}%%%%%%%%%%%%%%%
 \textsc{N.\ D.\ Baruah, B.\ C.\ Berndt, and H.\ H.\ Chan}, 
 Ramanujan's series for {$1/\pi$}: a survey, 
 \emph{Amer.\ Math.\ Monthly} {\bf 116(7)} (2009), 567--587. 

\bibitem{BerggrenBorweinBorwein2004}%%%%%%%%%%%%%%%
 \textsc{L.\ Berggren, J.\ Borwein, and P.\ Borwein}, 
 Pi: a source book, 
 Springer-Verlag, New York (2004). 

\bibitem{Berkovich2024}%%%%%%%%%%%%%%%
 \textsc{A.\ Berkovich}, 
 On the {$q$}-binomial identities involving the {L}egendre symbol modulo 3, 
 \emph{Discrete Math.} {\bf 347(3)} (2024), Paper No. 113824, 8. 

\bibitem{Berndt1994}%%%%%%%%%%%%%%%
 \textsc{B.\ C.\ Berndt}, 
 Ramanujan's notebooks. {P}art {IV}, 
 Springer-Verlag, New York (1994). 

\bibitem{BorweinBorwein1987}%%%%%%%%%%%%%%%
 \textsc{J.\ M.\ Borwein and P.\ B.\ Borwein}, 
 Pi and the {AGM}, 
 John Wiley \& Sons, Inc., New York (1987). 

\bibitem{Campbell2025further}%%%%%%%%%%%%%%%
 \textsc{J.\ M.\ Campbell}, 
 A further {$q$}-analogue of a formula due to {G}uillera, 
 \emph{Arch.\ Math.\ (Basel)} {\bf 124(5)} (2025), 511--516. 

 \bibitem{Campbell2025Hypergeometric}%%%%%%%%%%%%%%%
 \textsc{J.\ M.\ Campbell}, 
 Hypergeometric accelerations with shifted indices, 
 \emph{J.\ Difference Equ.\ Appl.} {\bf 31(9)} (2025), 1169--1192. 

\bibitem{Campbellunpublished}%%%%%%%%%%%%%%%
 \textsc{J.\ M.\ Campbell}, 
 An iterative approach toward hypergeometric accelerations, 
 arXiv:2505.14392 (2025). 
 
\bibitem{Campbell2026}%%%%%%%%%%%%%%%
 \textsc{J.\ M.\ Campbell}, 
 {$q$}-analogues of {\it {$\pi $}}-formulas due to {R}amanujan and {G}uillera, 
 \emph{J.\ Symbolic Comput.} {\bf 134} (2026), Paper No. 102531. 

\bibitem{CampbellLevrie2024}%%%%%%%%%%%%%%%
 \textsc{J.\ M.\ Campbell and P.\ Levrie}, 
 On two-term hypergeometric recursions with free lower parameters, 
 \emph{J.\ Difference Equ.\ Appl.} {\bf 30(7)} (2024), 870--891. 

\bibitem{Carmichael1912}%%%%%%%%%%%%%%%
 \textsc{R.\ D.\ Carmichael}, 
 The general theory of linear {{\(q\)}}-difference equations, 
 \emph{Am.\ J.\ Math.} {\bf 34} (1912), 147--168. 
 
\bibitem{ChenChu2021}%%%%%%%%%%%%%%%
 \textsc{X.\ Chen and W.\ Chu}, 
 Hidden {$q$}-analogues of {R}amanujan-like {$\pi$}-series, 
 \emph{Ramanujan J.} {\bf 54(3)} (2021), 625--648. 

\bibitem{ChenChu2023}%%%%%%%%%%%%%%%
 \textsc{X.\ Chen and W.\ Chu}, 
 {$q$}-analogues of {$\pi$}-series by applying {C}arlitz inversions to {$q$}-{P}faff-{S}aalsch\"utz theorem, 
 \emph{Contrib.\ Discrete Math.} {\bf 18(2)} (2023), 91--112. 

\bibitem{ChenChu2024}%%%%%%%%%%%%%%%
 \textsc{X.\ Chen and W.\ Chu}, 
 Triplicate dual series of the {D}ougall-{D}ixon theorem, 
 \emph{Rocky Mountain J.\ Math.} {\bf 54(2)} (2024), 407--422. 

\bibitem{Chu2011}%%%%%%%%%%%%%%%
 \textsc{W.\ Chu}, 
 Dougall's bilateral {${}_2H_2$}-series and {R}amanujan-like {$\pi$}-formulae, 
 \emph{Math.\ Comp.} {\bf 80(276)} (2011), 2223--2251. 

\bibitem{Chu2021GouldHsu}%%%%%%%%%%%%%%%
 \textsc{W.\ Chu}, 
 Ramanujan-like formulae for {$\pi$} and {$1/\pi$} via {G}ould-{H}su inverse series relations, 
 \emph{Ramanujan J.} {\bf 56(3)} (2021), 1007--1027. 

\bibitem{Chu2021Omega}%%%%%%%%%%%%%%%
 \textsc{W.\ Chu}, 
 Infinite series identities derived from the very well-poised {$\Omega$}-sum, 
 \emph{Ramanujan J.} {\bf 55(1)} (2021), 239--270. 

\bibitem{Chu1995}%%%%%%%%%%%%%%%
 \textsc{W.\ C.\ Chu}, 
 Inversion techniques and combinatorial identities---{J}ackson's {$q$}-analogue of the
 {D}ougall-{D}ixon theorem and the dual formulae, 
 \emph{Compositio Math.} {\bf 95(1)} (1995), 43--68. 

\bibitem{Chu2018}%%%%%%%%%%%%%%%
 \textsc{W.\ Chu}, 
 {$q$}-series reciprocities and further {$\pi$}-formulae, 
 \emph{Kodai Math. J.} {\bf 41(3)} (2018), 512--530. 

\bibitem{ChuZhang2014}%%%%%%%%%%%%%%%
 \textsc{W.\ Chu and W.\ Zhang}, 
 Accelerating {D}ougall's {${}_5F_4$}-sum and infinite series involving {$\pi$}, 
 \emph{Math.\ Comp.} {\bf 83(285)} (2014), 475--512. 

\bibitem{Fabry1911}%%%%%%%%%%%%%%%
 \textsc{E.\ Fabry}, 
 Th\'{e}orie des s\'{e}ries \`{a} terms constants: applications aux calculs num\'{e}riques, 
 Hermann, Paris (1911). 

\bibitem{GasperRahman2004}%%%%%%%%%%%%%%%
 \textsc{G.\ Gasper and M.\ Rahman}, 
 Basic hypergeometric series, 
 Cambridge University Press, Cambridge (2004). 
 
\bibitem{Guillera2006}%%%%%%%%%%%%%%%
 \textsc{J.\ Guillera}, 
 Generators of some {R}amanujan formulas, 
 \emph{Ramanujan J.} {\bf 11(1)} (2006), 41--48. 

\bibitem{Guillera2008}%%%%%%%%%%%%%%%
 \textsc{J.\ Guillera}, 
 Hypergeometric identities for 10 extended {R}amanujan-type series, 
 \emph{Ramanujan J.} {\bf 15(2)} (2008), 219--234. 

\bibitem{Guillera2002}%%%%%%%%%%%%%%%
 \textsc{J.\ Guillera}, 
 Some binomial series obtained by the {WZ}-method, 
 \emph{Adv. in Appl. Math.} {\bf 29(4)} (2002), 599--603. 

\bibitem{Guillera2018}%%%%%%%%%%%%%%%
 \textsc{J.\ Guillera}, 
 W{Z} pairs and {$q$}-analogues of {R}amanujan series for {$1/\pi$}, 
 \emph{J.\ Difference Equ.\ Appl.} {\bf 24(12)} (2018), 1871--1879. 

\bibitem{Guo2020}%%%%%%%%%%%%%%%
 \textsc{V.\ J.\ W.\ Guo}, 
 {$q$}-analogues of three {R}amanujan-type formulas for {$1/\pi$}, 
 \emph{Ramanujan J.} {\bf 52(1)} (2020), 123--132. 

\bibitem{GuoLiu2018}%%%%%%%%%%%%%%%
 \textsc{V.\ J.\ W.\ Guo and J.-C.\ Liu}, 
 {$q$}-analogues of two {R}amanujan-type formulas for {$1/\pi$}, 
 \emph{J. Difference Equ.\ Appl.} {\bf 24(8)} (2018), 1368--1373. 

\bibitem{GuoZudilin2018}%%%%%%%%%%%%%%%
 \textsc{V.\ J.\ W.\ Guo and W.\ Zudilin}, 
 Ramanujan-type formulae for {$1/\pi$}: {$q$}-analogues, 
 \emph{Integral Transforms Spec.\ Funct.} {\bf 29(7)} (2018), 505--513. 
 
\bibitem{HessamiPilehroodHessamiPilehrood2011}%%%%%%%%%%%%%%%
 \textsc{Kh.\ Hessami Pilehrood and T.\ Hessami Pilehrood}, 
 A {$q$}-analogue of the {B}ailey-{B}orwein-{B}radley identity, 
 \emph{J.\ Symbolic Comput.} {\bf 46(6)} (2011), 699--711. 

\bibitem{HouKrattenthalerSun2019}%%%%%%%%%%%%%%%
 \textsc{Q.-H.\ Hou, C.\ Krattenthaler, and Z.-W.\ Sun}, 
 On {$q$}-analogues of some series for {$\pi$} and {$\pi^2$}, 
 \emph{Proc.\ Amer.\ Math.\ Soc.} {\bf 147(5)} (2019), 1953--1961. 

\bibitem{HouSun2021}%%%%%%%%%%%%%%%
 \textsc{Q.-H.\ Hou and Z.-W.\ Sun}, 
 {$q$}-analogues of some series for powers of {$\pi$}, 
 \emph{Ann.\ Comb.} {\bf 25(1)} (2021), 167--177. 

\bibitem{Jackson1910q}%%%%%%%%%%%%%%%
 \textsc{F.\ H.\ Jackson}, 
 {{\(q\)}}-difference equations, 
 \emph{Am.\ J.\ Math.} {\bf 32} (1910), 305--314. 

\bibitem{Jackson1910Transformations}%%%%%%%%%%%%%%%
 \textsc{F.\ H.\ Jackson}, 
 Transformations of $q$-series, 
 \emph{Messenger of Math.} {\bf 39} (1910), 145--153. 

\bibitem{PetkovsekWilfZeilberger1996}%%%%%%%%%%%%%%%
 \textsc{M.\ Petkov\v sek, H.\ S.\ Wilf, and D.\ Zeilberger}, 
 {{$A=B$}}, 
 A K Peters, Ltd., Wellesley, MA (1996). 
 
\bibitem{vanderPoorten197879}%%%%%%%%%%%%%%%
 \textsc{A.\ van der Poorten}, 
 A proof that {E}uler missed{$\ldots $}Ap\'ery's proof of the irrationality of {$\zeta (3)$}, 
 \emph{Math.\ Intelligencer} {\bf 1(4)} (1978/79), 195--203. 

\bibitem{Rahman1993}%%%%%%%%%%%%%%%
 \textsc{M.\ Rahman}, 
 Some quadratic and cubic summation formulas for basic hypergeometric series, 
 \emph{Canad.\ J.\ Math.} {\bf 45(2)} (1993), 394--411. 

\bibitem{Rainville1960}%%%%%%%%%%%%%%%
 \textsc{E.\ D.\ Rainville}, 
 Special functions, 
 The Macmillan Company, New York (1960). 

 \bibitem{Ramanujan1914}%%%%%%%%%%%%%%%
 \textsc{S.\ Ramanujan}, 
 Modular equations and approximations to $\pi$, 
 \emph{Quart.\ J.} {\bf 45} (1914), 350--372. 

\bibitem{Starcher1931}%%%%%%%%%%%%%%%
 \textsc{G.\ W.\ Starcher}, 
 On {I}dentities {A}rising from {S}olutions of {$q$}-{D}ifference {E}quations and {S}ome {I}nterpretations in {N}umber {T}heory, 
 \emph{Amer.\ J.\ Math.} {\bf 53(4)} (1931), 801--816. 

\bibitem{Sun2019}%%%%%%%%%%%%%%%
 \textsc{Z.-W.\ Sun}, 
 Two {$q$}-analogues of {E}uler's formula {$\zeta(2)=\pi^2/6$}, 
 \emph{Colloq.\ Math.} {\bf 158(2)} (2019), 313--320. 

\bibitem{Trjitzinsky1933}%%%%%%%%%%%%%%%
 \textsc{W.\ J.\ Trjitzinsky}, 
 Analytic theory of linear {$q$}-difference equations, 
 \emph{Acta Math.} {\bf 61(1)} (1933), 1--38. 

 \bibitem{WangZhong2023}%%%%%%%%%%%%%%%
 \textsc{R.-H.\ Wang and M.\ X. X.\ Zhong}, 
 {$q$}-rational reduction and {$q$}-analogues of series for {$\pi$}, 
 \emph{J.\ Symbolic Comput.} {\bf 116} (2023), 58--71. 
 
\bibitem{Wei2023}%%%%%%%%%%%%%%%
 \textsc{C.\ Wei}, 
 On two double series for {$\pi$} and their {$q$}-analogues, 
 \emph{Ramanujan J.} {\bf 60(3)} (2023), 615--625. 

\bibitem{Wei2020}%%%%%%%%%%%%%%%
 \textsc{C.\ Wei}, 
 {$q$}-analogues of several {$\pi$}-formulas, 
 \emph{Proc.\ Amer.\ Math.\ Soc.} {\bf 148(6)} (2020), 2287--2296. 

\bibitem{WeiRuan2024}%%%%%%%%%%%%%%%
 \textsc{C.\ Wei and G.\ Ruan}, 
 Some double series for {$\pi$} and their {$q$}-analogues, 
 \emph{J.\ Math.\ Anal.\ Appl.} {\bf 537(1)} (2024), Paper No. 128309, 11. 

\bibitem{Wilf1999}%%%%%%%%%%%%%%%
 \textsc{H.\ S.\ Wilf}, 
 Accelerated series for universal constants, by the {WZ} method, 
 \emph{Discrete Math.\ Theor.\ Comput.\ Sci.} {\bf 3(4)} (1999), 189--192. 

\bibitem{Zeilberger1993}%%%%%%%%%%%%%%%
 \textsc{D.\ Zeilberger}, 
 Closed form (pun intended!), 
 \emph{A tribute to {E}mil {G}rosswald: number theory and related analysis} {\bf 143} (1993), 579--607, 
 Amer.\ Math.\ Soc., Providence, RI. 

\bibitem{Zeilberger1991}%%%%%%%%%%%%%%%
 \textsc{D.\ Zeilberger}, 
 The method of creative telescoping, 
 \emph{J.\ Symbolic Comput.} {\bf 11(3)} (1991), 195--204. 

\end{thebibliography}

 \ 

{\textsc{John M. Campbell}} 

\vspace{0.1in}

Department of Mathematics and Statistics

Dalhousie University 

6283 Alumni Crescent, Halifax, NS B3H 4R2

\vspace{0.1in}

{\tt jh241966@dal.ca}

\end{document}
